\documentclass[10pt,a4paper]{article}

% ----------------- Paquetes -------------------%
\usepackage[T1]{fontenc}
\usepackage[utf8]{inputenc}
\usepackage[a4paper,margin=2.5cm]{geometry}
\usepackage{mathptmx}             
\usepackage{changepage} 
\usepackage{setspace}
\usepackage[spanish]{babel}
\usepackage[labelfont=bf,labelsep=space]{caption}
\usepackage{amssymb}
\usepackage{amsmath}
\usepackage{ebgaramond}
\usepackage{placeins}
\usepackage{etoolbox}
\usepackage{setspace}
\usepackage{parskip} 
\usepackage{tcolorbox}
\usepackage{needspace}
\usepackage{xparse}
\usepackage{ocgx2}
\usepackage{hyperref}
\usepackage{hypcap}


%%%%%%%%%%%%%%%%%%%%%%%%%%%%%%%%%%%%%%%%%%%%%%%%%%%%%%%%%%%%%%%%
% --- Fija primero tus valores globales "buenos" ---
\setstretch{1.2}                 % Interlineado global
\setlength{\parskip}{0.7\baselineskip}
\setlength{\parindent}{0pt}
\newlength{\GlobalParskip}
\setlength{\GlobalParskip}{\parskip}

\newcounter{eqnum}[subsection]
\renewcommand{\theeqnum}{\arabic{eqnum}}
\newcounter{alphaitem}
\newcounter{numitem}

%% ------------------- Recuadros----------------------%%
\newcommand{\puntosclave}[1]{%
  \begin{tcolorbox}[
    colframe=black,
    title={Puntos clave},
    before upper={\setlength{\parindent}{0.5cm}\setlength{\parskip}{0.5\baselineskip}}
  ]
  #1
  \end{tcolorbox}
}

\newcommand{\modificaciones}[1]{
  \begin{tcolorbox}[
    colback=white,
    colframe=red,
    title={Modificaciones a realizar},
    before upper={\setlength{\parindent}{0.5cm}\setlength{\parskip}{0.5\baselineskip}}
  ]
  #1
  \end{tcolorbox}
}

%% ------------------- Listas----------------------%%
    % --- Lista alfabética ---
    \newenvironment{la}
      {\begin{list}{\alph{alphaitem})}{%
          \usecounter{alphaitem}%
          \setlength{\leftmargin}{1cm}%
          \setlength{\parskip}{3pt}% 
          \setlength{\itemsep}{0pt}% ← espacio entre ítems
          \setlength{\parsep}{0pt}%  ← párrafos dentro de un ítem
      }}
      {\end{list}}
      
    % --- Lista numérica ---
    \newenvironment{lnum}
      {\begin{list}{\arabic{numitem}.}{%
          \usecounter{numitem}%
          \setlength{\leftmargin}{1cm}%
          \setlength{\parskip}{3pt}% 
          \setlength{\itemsep}{0pt}% ← espacio entre ítems
          \setlength{\parsep}{0pt}%  ← párrafos dentro de un ítem
      }}
      {\end{list}}
      
% ------------------- Imágenes----------------------%
    \graphicspath{{Img/}}% Rutas por defecto 
    % --- Imagen adaptativa ---
    % Registros de dimensión definidos una sola vez
    \newdimen\imgcap
    \newdimen\imgmin
    \newdimen\imgmax
    \newdimen\imgremain
    \newdimen\imgh
    \newdimen\imgneed
    
    \newcommand{\imagenfit}[3]{%
      \addvspace{10pt}% separación antes
      \par\begingroup
        % Parámetros de composición (asignaciones, no \newdimen)
        \imgcap=1\baselineskip            % alto estimado de título+fuente
        \imgmin=0.15\textheight
        \imgmax=0.15\textheight
        \imgremain=\pagegoal
        \ifdim\pagegoal=\maxdimen \imgremain=0pt\fi
        \advance\imgremain by -\pagetotal
        \advance\imgremain by -\imgcap
        % Decisión de altura destino
        \ifdim\imgremain<\imgmin
          \newpage
          \imgh=0.20\textheight
        \else
          \imgh=\imgremain
          \ifdim\imgh>\imgmax \imgh=\imgmax \fi
        \fi
        % Reservar espacio para evitar cortes del bloque completo
        \imgneed=\imgh \advance\imgneed by \imgcap
        \Needspace{\imgneed}
        % Cuerpo
        \noindent\begin{minipage}{\textwidth}
          \centering
          \captionof{figure}{\textemdash\ #2}\par\vspace{0.3em}% título arriba
          \includegraphics[width=\linewidth,height=\imgh,keepaspectratio]{\detokenize{#1}}\par
          \vspace{0.3em}\small\textit{Fuente: }#3% fuente abajo
        \end{minipage}%
      \endgroup
      \par\addvspace{10pt}%
    }

%------------ Bloque de RECUADRO ESPECIAL -------------%
    \newenvironment{specialblock}[1]{%
      \begin{quote} % crea sangría a izquierda y derecha
      \small % texto más pequeño
      \noindent\textbf{#1}\\[-0.3em] % título en negrita
      \rule{\linewidth}{0.4pt}\\[0.6em]}{
      \end{quote}}
      

%-------------------------- Author Font -----------------------%
    \newcommand{\authorfont}[1]{{\fontfamily{EBGaramond-TLF}\selectfont #1}}

%-------------------------- Anexos -----------------------%
    \makeatletter
    \newcounter{anexo}
    \renewcommand{\theanexo}{\Roman{anexo}}
    \newcommand{\anexo}[1]{%
      \clearpage
      \phantomsection
      \refstepcounter{anexo}%
      \noindent\textbf{Anexo \theanexo.- }#1\par\medskip
      \addcontentsline{stc}{section}{Anexo \theanexo.- #1}%
    }
    \makeatother

    %-------------------------- Lista de figuras (personal) -----------------------%
\makeatletter
\newcommand{\MyListOfFigures}{%
  \clearpage
  \phantomsection
  {\centering\bfseries\Large Índice de figuras\par}%
  \medskip
  % Entrada en nuestro índice de contenido (stc)
  \addcontentsline{stc}{section}{Índice de figuras}%
  % Recuperación del listado (sin cabecera propia)
  \begingroup
    \parindent=0pt\parskip=2pt
    \@starttoc{lof}%
  \endgroup
}
\makeatother

%-------------------------- Bibliografía -----------------------%
\makeatletter
% Contador para las referencias
\newcounter{myref}
% Cabecera de la bibliografía, entra en el índice 'stc'
\newcommand{\MyBibliography}{%
  \clearpage
  \phantomsection
  {\centering\bfseries\Large Bibliografía y referencias\par}%
  \medskip
  \addcontentsline{stc}{section}{Bibliografía y referencias}%
  \setcounter{myref}{0}%
}
\newcommand{\MyBibItem}[4]{%
  \refstepcounter{myref}%
  \noindent[\themyref]\ #1.\ \emph{#2}. #3. #4\par
}
\makeatother

%--------- Establecer cuatro niveles de índice --------%
    \setcounter{tocdepth}{4}   
    \setcounter{secnumdepth}{4}% numera hasta \paragraph

%%%%%%%%%%%%%%%%%%%%%%%%%%%%%%%%

\makeatletter

    %--------- Interlineado Global --------%
    \edef\GlobalBaseLineStretch{\baselinestretch}
    
    %--------- Espaciado de seccinones --------%
    \renewcommand\section{\@startsection{section}
      {1}{\z@}
      {2.5ex \@plus 1ex \@minus .2ex} % espacio antes
      {1.5ex \@plus .2ex}             % espacio después
      {\normalfont\Large\bfseries}    % formato del título
    }
    \renewcommand\subsection{\@startsection{subsection}{2}{\z@}
      {3ex \@plus 0.8ex \@minus .2ex}
      {1.5ex \@plus .2ex}
      {\normalfont\large\bfseries}
    }
    \renewcommand\subsubsection{\@startsection{subsubsection}{3}{\z@}
      {3ex \@plus 0.5ex \@minus .2ex}
      {1ex \@plus .2ex}
      {\normalfont\normalsize\bfseries}
    }
    \renewcommand\paragraph{\@startsection{paragraph}{4}{\z@}%
    {-2ex \@plus -0.5ex \@minus -.2ex}% espacio antes
    {0.8ex \@plus .2ex}% espacio después
    {\normalfont\normalsize\itshape}}% ← cursiva en lugar de negrita

    %--------- Ecuaciones --------%   
    % Variables globales para almacenar títulos y notas
    \gdef\leq@current@section{}%
    \gdef\leq@current@subsection{}%
    \gdef\leq@last@printed@section{}%
    \gdef\leq@last@printed@subsection{}%
    
    % Comando para añadir nota (invisible en el cuerpo)
    \newcommand{\eqnote}[1]{%
      \addcontentsline{leq}{leqnote}{#1}%
    }
    
    % Comando para ecuaciones
    \newcommand{\eqblock}[2]{%
      % Verificar si necesitamos imprimir el título de sección
      \ifx\leq@current@section\leq@last@printed@section\else
        \addcontentsline{leq}{leqsection}{\leq@current@section}%
        \global\let\leq@last@printed@section\leq@current@section
        \gdef\leq@last@printed@subsection{}% Reset subsección al cambiar sección
      \fi
      % Verificar si necesitamos imprimir el título de subsección
      \ifx\leq@current@subsection\leq@last@printed@subsection\else
        \ifx\leq@current@subsection\@empty\else
          \addcontentsline{leq}{leqsubsection}{\leq@current@subsection}%
          \global\let\leq@last@printed@subsection\leq@current@subsection
        \fi
      \fi
      % Ecuación
      \refstepcounter{eqnum}%
      \begin{equation*}
        #1\tag{E.\theeqnum}
      \end{equation*}%
      \addcontentsline{leq}{leqt}{[E.\theeqnum] -- #2}%
      \addcontentsline{leq}{leqm}{#1}%
    }
    
    %%% ----- Índice de ecuaciones ----------%%%
    % Tamaño para []
    \newcommand{\leqIndexFont}{\normalsize}
    \newcommand{\leqTitleFont}{\tiny}
    \newcommand{\leqNoteFont}{\tiny}
    % Cabecera de sección 
    \newcommand*\l@leqsection[2]{%
      \addvspace{25pt}% Mayor separación antes de nueva sección
      \noindent\leqIndexFont
      \makebox[\linewidth]{\rule{.38\linewidth}{0.4pt}\ \ \textbf{  \MakeUppercase{#1}}\ \ \rule{.38\linewidth}{0.4pt}}%
      \par\vspace{-10pt}
    }
    % Cabecera de subsección 
    \newcommand*\l@leqsubsection[2]{%
      \par\addvspace{15pt}%
      \noindent\leqIndexFont #1
      \par\vspace{-7pt}
    }
    % Título de ecuación 
    \newcommand*\l@leqt[2]{%
    \par\addvspace{10pt}%
      \par\noindent\leqTitleFont #1\par
    }
    % Miniatura de ecuación
    \newcommand*\l@leqm[2]{%
      \vspace{0.3\baselineskip}%
      \par\scriptsize\noindent\hspace*{1.5em}\(#1\)\par%
    }
    % Nota de ecuación 
    \newcommand*\l@leqnote[2]{%
      \vspace{0.2\baselineskip}%
      \par\noindent\leqNoteFont\hspace*{1.5em}\textbf{Nota.--} #1\par\vspace{0.5\baselineskip}%
    }
    % Generación del índice
    \newcommand{\mylistofequations}{%
      \clearpage
      \phantomsection
      % Título visual de la lista (ajústalo a tu gusto):
      {\centering\bfseries\Large Índice de ecuaciones\par}
      % Entrada en nuestro índice 'stc':
      \addcontentsline{stc}{section}{Índice de ecuaciones}%
      % Llamada a tu lista real (usa el nombre exacto de tu macro):
      \@starttoc{leq}%
    }
    
    %--------- Captura de títulos de sección --------%
    \let\leq@original@section\section
    \let\leq@original@subsection\subsection
    % Flag para saber si estamos en una sección asterisco
    \newif\ifLeqStarSection
    \LeqStarSectionfalse
    
    % Redefinir \section CON CONTROL DE ASTERISCO
    \renewcommand{\section}{%
      \@ifstar{\leq@section@star}{\@dblarg\leq@section@nostar}%
    }
    \def\leq@section@star#1{%
      \global\LeqStarSectiontrue
      \gdef\leq@current@section{#1}%
      \gdef\leq@current@subsection{}%
      \leq@original@section*{#1}%
      \global\LeqStarSectionfalse
    }
    \def\leq@section@nostar[#1]#2{%
      \global\LeqStarSectionfalse
      \gdef\leq@current@section{#2}%
      \gdef\leq@current@subsection{}%
      \leq@original@section[#1]{#2}%
    }
    % Redefinir \subsection
    \renewcommand{\subsection}{%
      \@ifstar{\leq@subsection@star}{\@dblarg\leq@subsection@nostar}%
    }
    \def\leq@subsection@star#1{%
      \global\LeqStarSectiontrue
      \gdef\leq@current@subsection{#1}%
      \leq@original@subsection*{#1}%
      \global\LeqStarSectionfalse
    }
    \def\leq@subsection@nostar[#1]#2{%
      \global\LeqStarSectionfalse
      \gdef\leq@current@subsection{#2}%
      \leq@original@subsection[#1]{#2}%
    }

    % ----- Restauración de valores Globales de estilo ----------%
    % Flags
    \newif\ifInFootnote
    \InFootnotefalse
    \pretocmd{\@footnotetext}{\InFootnotetrue}{}{}
    \apptocmd{\@footnotetext}{\InFootnotefalse}{}{}
    % Profundidad de listas (soporta anidadas)
    \newcount\MyListDepth \MyListDepth=0
    \AtBeginEnvironment{la}{\advance\MyListDepth by 1}
    \AfterEndEnvironment{la}{\advance\MyListDepth by -1}
    \AtBeginEnvironment{lnum}{\advance\MyListDepth by 1}
    \AfterEndEnvironment{lnum}{\advance\MyListDepth by -1}
    \AtBeginEnvironment{itemize}{\advance\MyListDepth by 1}
    \AfterEndEnvironment{itemize}{\advance\MyListDepth by -1}
    \AtBeginEnvironment{enumerate}{\advance\MyListDepth by 1}
    \AfterEndEnvironment{enumerate}{\advance\MyListDepth by -1}
    % Restauración diferida: hazlo al INICIO del próximo párrafo real
    \newif\if@pendingRestore
    \@pendingRestorefalse
    \newcommand{\RequestRestore}{\global\@pendingRestoretrue}
    \everypar{%
      \if@pendingRestore
        \global\@pendingRestorefalse
        \ifInFootnote\else
          \ifnum\MyListDepth=0
            \setlength{\parskip}{\GlobalParskip}%
            \setstretch{\GlobalBaseLineStretch}%
          \fi
        \fi
      \fi
    }
    % Al final de bloques:
    \AfterEndEnvironment{equation}{\RequestRestore}
    \AfterEndEnvironment{equation*}{\RequestRestore}
    \AfterEndEnvironment{la}{\RequestRestore}
    \AfterEndEnvironment{lnum}{\RequestRestore}
    % Al final de macros:
    \apptocmd{\eqblock}{\RequestRestore}{}{}
    \apptocmd{\imagenfit}{\RequestRestore}{}{}
    
\makeatother

% ===================== ToC propio con 4 niveles (único bloque) =====================
\makeatletter

% --- Escritores: añadimos una línea al .stc en cada cabecera numerada ---
% Guardamos las versiones actuales para llamar al comportamiento normal
\let\MyTOC@orig@section\section
\let\MyTOC@orig@subsection\subsection
\let\MyTOC@orig@subsubsection\subsubsection
\let\MyTOC@orig@paragraph\paragraph

% SECTION
\renewcommand{\section}{%
  \@ifstar{\MyTOC@section@star}{\@dblarg\MyTOC@section@nostar}%
}
\def\MyTOC@section@star#1{%
  \MyTOC@orig@section*{#1}% sin entrada al .stc
}
\def\MyTOC@section@nostar[#1]#2{%
  \MyTOC@orig@section[{#1}]{#2}% comportamiento normal (escribe al .toc)
  % Escribimos también nuestra línea al .stc (texto con \numberline)
  \addcontentsline{stc}{section}{\protect\numberline{\thesection}#2}%
}

% SUBSECTION
\renewcommand{\subsection}{%
  \@ifstar{\MyTOC@subsection@star}{\@dblarg\MyTOC@subsection@nostar}%
}
\def\MyTOC@subsection@star#1{%
  \MyTOC@orig@subsection*{#1}%
}
\def\MyTOC@subsection@nostar[#1]#2{%
  \MyTOC@orig@subsection[{#1}]{#2}%
  \addcontentsline{stc}{subsection}{\protect\numberline{\thesubsection}#2}%
}

% SUBSUBSECTION
\renewcommand{\subsubsection}{%
  \@ifstar{\MyTOC@subsubsection@star}{\@dblarg\MyTOC@subsubsection@nostar}%
}
\def\MyTOC@subsubsection@star#1{%
  \MyTOC@orig@subsubsection*{#1}%
}
\def\MyTOC@subsubsection@nostar[#1]#2{%
  \MyTOC@orig@subsubsection[{#1}]{#2}%
  \addcontentsline{stc}{subsubsection}{\protect\numberline{\thesubsubsection}#2}%
}

% PARAGRAPH
\renewcommand{\paragraph}{%
  \@ifstar{\MyTOC@paragraph@star}{\@dblarg\MyTOC@paragraph@nostar}%
}
\def\MyTOC@paragraph@star#1{%
  \MyTOC@orig@paragraph*{#1}%
}
\def\MyTOC@paragraph@nostar[#1]#2{%
  \MyTOC@orig@paragraph[{#1}]{#2}%
  % Si no numeras los \paragraph, cambia \theparagraph por texto plano sin \numberline
  \addcontentsline{stc}{paragraph}{\protect\numberline{\theparagraph}#2}%
}

% --- Impresor del índice propio (lee el .stc) ---
\newcommand{\MyTOC}{%
  \par\bigskip
  {\centering\bfseries\Large Índice de contenido\par}%
  \medskip
  \begingroup
    \parindent=0pt\parskip=2pt
    \@starttoc{stc}%
  \endgroup
}

\makeatother
% ===================== fin ToC propio =====================


%%%%%%%%%%%%%%


% --- Datos del tema ---
\title{Sistema Fiscal en España (I): Impuesto sobre la Renta de las Personas Físicas\\
\large }
\author{Marcelo Ortega}
\date{Septiembre 2026}

\begin{document}
\maketitle
\newpage

% --- Página de resumen y modificaciones ---
\puntosclave{ %% Escribir puntos clave Apuntes CECO
     Para facilidad del opositor, se ha optado por establecer no sólo una introducción común a los cuatro temas. sino además una estructura paralela entre los dos temas de imposición directa (éste yel3Bl7). Son necesarios los siguientes apuntes:
     \begin{la}
         \item La aproximación empleada es eminentemente jurídica. Se considera importante mantener la corrección formal para presentar correctamente estos temas. Sin embargo, la exposición se enriquecería notablemente si los temas 4B 16, 4B 17 y 4B 19 incluyen referencias laterales a las nociones estudiadas en el tema 4B 12 ( de manera similar para el tema 4B20 respecto del 4B 13); y,
         \item Año a afio se introducen pequeilas modificaciones en los impuestos. Es necesario por tanto actualizar el contenido del tema. pero sin perder de vista que los pequeilos detalles no son lo que aportarán valor en la exposición, sino el ser capaz de demostrar al tribunal que se comprenden las principales características del disefio de cada impuesto. así como su motivación.
     \end{la}

     \textcolor{red}{\textbf{OJO -> Extensiones interesantes:} Anexo (I), Determinación de la unidad de gravamen (tributación familiar); Anexo (II), Regímenes especiales del IRPF}

     \textcolor{red}{\textbf{OJO ->} Revisar la siguiente propuesta:} Esta no es una oposición jurídica, es una oposición cuyo objeto de estudio se encuentra en el puntto medio entre la academia y la aplicación técnica de los conocimientos por ésta acumulados a lo largo de los años. El Sistema Fiscal español, por tanto, puede abordarse desde dos perspectivas: (i) la presente en los temas de ICEX-CECO, donde se pone el énfasis en la estructura jurídica del impuesto, manteniendo un esquema lineal en el que presentamos el impuesto, luego su esquema de determinación y acabamos una síntesis; o, (ii) desde una óptica de Teoría Económica, en la que podemos empezar presentando el marco jurídico-conceptual del presente impuesto (historia, sujeto pasivo, régimen de gestión y liquidación, etc.), para luego presentar (en dos apartados) los objetos de estudio económico: la determinación de la renta sujeta a gravamen y los gastos fiscales que contempla la normativa del impuesto (deducciones y exenciones)
    }
\vspace{1cm}

\modificaciones{
    \textbf{NOTA (04/10/2026):} Revisar: https://x.com/rdomenechv/status/2047185603372913117?s=20

    \textbf{NOTA (04/10/2026):} Revisar: https://x.com/JuanLuis\_JG/status/2047193416337969188?s=20
}

\newpage
\MyTOC
\newpage

\mylistofequations
\clearpage

\MyListOfFigures
\clearpage


\pagenumbering{arabic}
\setcounter{page}{1}


\section{Introducción}
\textcolor{\textbf{REVISAR:} Se propone como introducción común para todos los temas de Sistema Fiscal español (16-20)}

La principal fuente de ingresos públicos está constituida por los ingresos tributarios, que gravan una determinada manifestación de capacidad económica, principio material básico del sistema tributarios español de acuerdo con el art.31. I de la Constitución Española. 

Los tributos pueden definirse, en consonancia con nuestra Ley General Tributaria, como los ingresos públicos que consisten en prestaciones pecuniaria5 exigidas por una Administración pública como consecuencia de la realización del supuesto de hecho al que la ley vincula el deber de contribuir, con el fin primordial de obtener los ingresos necesarios para el sostenimiento de los gastos públicos.

Existen distintas calificaciones de tributos. Así, según los sujetos activos que los establecen y exigen, pueden ser, municipales. provinciales. autonómicos y estatales; según que su exacción siga o no el procedimiento legalmente establecido, se distingue entre tributos fiscales y parafiscales; según los criterios utilizados en su cuantificación. tributos fijos y variables: y por último. quizá la calificación más utilizada, es la que diferencia los tributos según sean impuestos. tasas o contribuciones especiales.

Centrándonos en esta última calificación, recogida en el ait.2 de la Ley General Tributaria, vamos a definir cada una de estas fi¡,'Uras tributarias (impuestos, tasas y contribuciones especiales) destacando las similitudes y diferencias entre cada una de ellas.

• Los impuestos son aquellos tributos exigidos sin contraprestación, cuyo hecho imponible· está constituido por negocios, actos o hechos que ponen de manifiesto la capacidad contributiva del contribuyente. A diferencia de las otras especies tributarias, no aparece contemplado en su ámbito objetivo ninguna actividad administrativa.

• Las lasas. sin embargo, son los tributos cuyo hecho imponible consiste en la utilización privativa o el aprovechamiento especial del dominio público, la prestación de servicios o la realización de actividades en régimen de Derecho Público que se refieran. afecten o beneficien, de modo pa1ticular, al obligado tributa1io. cuando los servicios o actividades no sean de solicitud o recepción voluntaria para los obligados tributarios o no se presten o realicen por el sector privado (imaginemos una tasa municipal solicitada por un bar para el establecimiento de una terraza en la vía pública).

En este punto merece hacer una mención especial a la diferencia de las tasas con los precios públicos, que son contraprestaciones pecuniarias que se satisfacen por la prestación de servicios o la realización de actividades efectuadas en régimen de Derecho público cuando, prestándose también tales servicios o actividades por el sector privado. sean de solicitud voluntaria por parte de los administrados (como puede ser un precio público por la entrada a una piscina municipal).

• Por último, resaltar la figura de las contribuciones especiales. como los tributos cuyo hecho imponible consiste en la obtención por el obligado tributario de un beneficio o de un aumento de valor de sus bienes, cdmo consecuencia de la realización de obras pública<; o del establecimiento o ampliación de servicios públicos (como ejemplo pensemos en la instalación de una red de distribución de aguas en un determinado barrio o calle). Vemos como en la contribución hay siempre una actividad administrativa. al contrario que en el impuesto; y, a diferencia de la tasa, aun cuando en las dos categorías debe concurrir una determinada actividad administrativa. la que da lugar al pago de la contribución especial está básicamente encaminada a la satisfacción de un interés general. mientras que en el caso de la tasa, la actividad administrativa está sustancialmente motivada por el particular y persigue la solución de problemas individuales.

Por lo que respecta a la figura de los impuestos en particular. caben varias calificaciones. Según su presupuesto o��jetivo. distinguimos entre impuestos personales, que son aquellos que no pueden concebirse sino en relación con una persona (Ej: IRPF). e impuestos reales, asentados sobre un presupuesto objetivo cuya naturaleza se determina con independencia del elemento personal (Ej: 181): según que se consideren o no las circunstancias personales del sujeto pasivo. hablamos de impuestos subjetivos (Ej: 1 RPF o IS) o impuestos objetivos (Ej: 1 V A): i;:n función del elemento temporal, impuestos periódicos (Ej: IRPF. IBI. IS) e impuestos instantáneos (Ej: IVA. ITP y A.ID): y, para terminar. según que la norma jurídica tributaria establezca o no la obligación de pago del impuesto a cargo del sujeto pasivo o de otra persona a través del mecanismo de repercusión. impositivos directos (Ej: l RPF) e impuestos indirectos (Ej: IVA).

Una vez centrados los impuestos dentro del sistema financiero y tributario español. a lo largo de varios temas. se va a ir haciendo referencia y explicando en mayor profundidad una selección de las principales figuras impositivas. En este sentido. se hará referencia al Impuesto sobre la Renta de las Personas Fisicas y de los No Residentes. el Impuesto sobre Sociedades. el Impuesto sobre el Valor Añadido e Impuestos Especiales y una serie de impuestos patrimoniales, de distintos ámbitos territoriales. como el Impuesto sobre el Patrimonio. el Impuesto sobre Sucesiones y Donaciones, el Impuesto sobre Transmisiones Patrimoniales y Actos Jurídicos Documentados o el Impuesto sobre Bienes Inmuebles.

\subsection{Contextualización}


\subsubsection{Relevancia}




\subsubsection{Problemática}



\subsection{Estructura de la exposición}




\clearpage
\section{Impuesto sobre la Renta de las Personas Físicas (IRPF): Marco jurídico-conceptual}
\puntosclave{
    El IRPF es un tributo de carácter personal y directo que grava la renta de las personas físicas: (i) sus rendimientos del trabajo; (ii) sus rendimientos del capital; (iii) sus rendi mientos de actividades económicas; (iv) sus ganancias y pérdidas patrimoniales; y, (v) las imputaciones de renta que se establezcan por la ley. con independencia del lugar donde se hubiesen producido y cualquiera que sea la residencia  del pagador.
}


\subsection{Desarrollo histórico}
El IRPF actual no es el resultado de una evolución lineal, sino de sucesivos modelos de imposición sobre la renta, cada uno respondiendo a un problema económico o a una exigencia constitucional concreta. Repasemos brevemente su recorrido histórico para comprender la arquitectura actual del impuesto.

\subsubsection{Antecedentes: el sistema de imposición de producto (hasta 1978)}
Antes de 1978, España carecía de un impuesto personal y sintético sobre la renta comparable a los de nuestro entorno. El sistema se articulaba mediante figuras de producto o cedulares que gravaban cada fuente de renta de forma aislada, sin acumulación ni progresividad efectiva sobre la renta total del contribuyente, y sin atender a sus circunstancias personales y familiares (Contribución Territorial (rústica y urbana), Contribución Industrial y de Comercio, Impuesto sobre los Rendimientos del Trabajo Personal, Impuesto sobre las Rentas del Capital). 

De hecho, el intento más temprano de introducir un impuesto general y personal, Contribución General sobre la Renta (1932/1933), tuvo una aplicación testimonial y no llegó a sustituir al sistema de producto. Es importante comprender este punto de partida porque permite entender por qué la reforma de 1978 no fue un simple cambio normativo, sino un cambio de paradigma técnico: de la imposición real y cedular a la imposición personal y sintética.

\subsubsection{Reforma Fuentes Quintana: Ley 44/1978}
La Ley 44/1978 del IRPF supuso la sustitución del modelo de tributación imperante hasta entonces mediante el establecimiento de un impuesto directo y personal que pretendía gravar con carácter sintético la totalidad de las rentas obtenidas por los sujetos pasivos, en función de su cuantía y de las circunstancias personales y familiares concurrentes. Por su encuadre dentro del proceso de transición, su aprobación se enmarca en la reforma tributaria diseñada por FUENTES QUINTANA como parte de los Pactos de la Moncloa (1977), en un contexto de crisis económica (inflación elevada, déficit fiscal estructural, sistema fiscal débil). Se aprueba también, como vemos, meses antes de la Constitución de 1978, cuyo art. $31.1$ CE consagrará después los principios de generalidad, igualdad, progresividad y capacidad económica que la reforma ya anticipaba materialmente.

Características técnicas reseñables son diversas, aunque destacan: (i) \textbf{impuesto sintético}, integración de todas las fuentes de renta en una base única, frente al modelo cedular anterior; (ii) configuración de la \textbf{unidad familiar} como sujeto pasivo, con \textbf{acumulación obligatoria} de las rentas de sus miembros; (iii) \textbf{tarifa progresiva única}. 

Este segundo rasgo, tributación familiar obligatoria, será precisamente el punto de fractura constitucional que da paso al siguiente hito.

\paragraph{Crisis constitucional de la unidad familiar: STC 45/1989 y Ley 20/1989}
La STC 45/1989, de 20 de febrero, declaró inconstitucional la acumulación obligatoria de rentas en el seno de la unidad familiar, obligando a configurar el IRPF como una exacción sobre la renta individual, lo que se hizo mediante la Ley 20/1989, de 28 de julio, de adaptación del IRPF y del Impuesto Extraordinario sobre el Patrimonio. 

El fundamento de fondo del Tribunal no fue prohibir la tributación conjunta como opción, sino declarar inconstitucional que la mera pertenencia a una unidad familiar determinase, por el efecto de la progresividad sobre una base acumulada, una carga tributaria superior a la que soportarían los mismos individuos tributando por separado, lo que constituye un ejemplo claro del famoso \textit{impuesto sobre el matrimonio} ([\authorfont{Ver Tema 4.B.12}]) y vulneraba los principios de capacidad económica e igualdad del art. 31.1 CE en relación con el art. 39 CE (protección de la familia y equidad horizontal entre familias).\footnote{\textbf{NOTA.-} La reforma sobre tributación conjunta/unidad familiar (STC 45/1989) vínculo el único episodio de la historia del IRPF español con rango de declaración de inconstitucionalidad.
}

La Ley 20/1989, por tanto, fue una adaptación urgente; la reforma sistemática y definitiva llegó después. 

\subsubsection{Ley 18/1991: Primer IRPF post-constitucional}
La Ley 18/1991, de 6 de junio, supone la primera reforma integral tras la doctrina constitucional, consolidando la tributación individual como regla y la conjunta como opción voluntaria (germen de los actuales arts. $82-84$ LIRPF), además de introducir mejoras técnicas en la determinación de rendimientos y en el tratamiento de las ganancias patrimoniales. 

Este será el modelo seguido durante buena parte de los años noventa, en un contexto de consolidación fiscal previo a la entrada en la Unión Económica y Monetaria.

\paragraph{Ley 40/1998: Reducción de tipos y la simplificación}
El Ministerio de Hacienda creó, el 17 de febrero de 1997, por Resolución de la Secretaría de Estado de Hacienda, una Comisión de expertos presidida por el catedrático LAGARES para el estudio y propuesta de medidas para la reforma del IRPF, cuyos trabajos se publicaron en febrero de 1998 y desembocaron, con una tramitación parlamentaria inusualmente rápida, en la Ley 40/1998, de 9 de diciembre, del IRPF, en vigor desde el 1 de enero de 1999. 

Los tres ejes declarados de la reforma eran la reducción generalizada de la carga tributaria (con especial atención a rentas medias y bajas), el fomento del ahorro privado y la inversión, y la simplificación del diseño y gestión del impuesto. Técnicamente, sustituye el sistema de deducciones en la cuota por circunstancias personales y familiares por el mecanismo del mínimo personal y familiar (renta exenta), germen conceptual que sigue vigente y forma parte de un debate académico más amplío que trataremos después: es la primera vez que el legislador articula explícitamente un mínimo exento con fundamento en capacidad económica, no como beneficio fiscal discrecional. 

Poco después, en 2004, se aprueba el Texto Refundido de la Ley del IRPF, sin alterar sustancialmente el modelo de 1998, con función de consolidación normativa. 

\subsubsection{Ley 35/2006: la dualización de fuentes de renta}
El verdadero salto de modelo llega con la Ley 35/2006, de 28 de noviembre, actualmente vigente (con numerosas modificaciones posteriores). 

Su rasgo definitorio es la separación entre renta general y renta del ahorro, con tarifas y reglas de integración diferenciadas: se abandona el modelo sintético puro de 1978 para aproximarse a un modelo de dualización inspirado en los sistemas nórdicos (\textit{Dual Income Tax} de Dinamarca, Suecia, Noruega, Finlandia desde finales de los ochenta), motivado por la creciente movilidad internacional del capital tras la liberalización de movimientos de capitales en la Unión: gravar el capital al mismo tipo marginal progresivo que el trabajo generaba fuertes incentivos a la elusión y a la deslocalización del ahorro, de modo que se opta por un tipo proporcional para la base del ahorro (inicialmente único, luego progresivo por tramos). 

Este es un punto en el que incidiremos más adelante: la dualización se corresponde con la primera gran ruptura deliberada y estructural respecto al benchmark de renta gravable como incremento neto de patrimonio gravado uniformemente, con independencia de su fuente (HAIG-SIMONS).

\paragraph{Crisis económica y el gravamen complementario (2012-2014)}
La crisis económica posterior a la reforma de 2006 redujo de forma significativa los ingresos tributarios, elevando el déficit público hasta niveles desconocidos, lo que llevó a aprobar un gravamen complementario en el IRPF, de aplicación en los ejercicios 2012, 2013 y 2014, con una elevación progresiva de la carga tributaria en función del nivel de renta. 

Este episodio ejemplifica muy bien el uso del IRPF como instrumento de consolidación: se puede subir la tarifa en caliente sin tocar el diseño legislativo, por ejemplo del hecho imponible. 

\paragraph{Ley 26/2014: el segundo Informe Lagares}
Durante el periodo álgido de la crisis, el Consejo de Ministros constituyó por acuerdo en 2013 una comisión de expertos para la reforma del sistema tributario, presidida de nuevo por LAGARES e integrada por ocho vocales más, análoga a la de 1997, que entregó su informe en 2014, con $125$ propuestas de reforma y $270$ modificaciones impositivas concretas. 

Entre sus críticas al IRPF vigente destacaba, según recoge el propio informe, una cierta confusión terminológica en la Ley 35/2006 entre renta y rendimiento que dificultaba la comprensión del impuesto por el contribuyente medio, además de proponer la recomposición de bases mediante la eliminación de exenciones, reducciones y bonificaciones dispersas (es decir, una crítica directa a la proliferación de gastos fiscales, que veremos en el último bloque). 

De ese informe deriva la Ley 26/2014, que modifica sustancialmente la Ley 35/2006 vigente. Los cambios más destacados de esta reforma fueron la reducción del número de tramos y de los tipos marginales, supresión de la exención parcial de dividendos, endurecimiento del régimen de exención por reinversión en vivienda habitual (limitado a adquisiciones anteriores a 2013), un nuevo tratamiento de la compensación de rendimientos negativos del ahorro con el general (hasta un límite del $25\%$, luego ampliado).

\paragraph{Reformas recientes: 2017-2026}
Los últimos años muestran un patrón distinto al de las grandes reformas integrales, ya que se han llevado a cabo muchos ajustes puntuales, pero casi siempre vía LPGE y dirigidos a incrementar la progresividad en la parte alta de la tarifa. 

La LPGE 11/2020, por ejemplo, añadió un nuevo tramo del $24,50\%$ a la base liquidable general que exceda de $300.000€$, y un nuevo tramo del $13\%$ a la base liquidable del ahorro que exceda de $200.000€$, con efectos desde el 1 de enero de 2021. Más tarde, la Ley 7/2024, incrementó en dos puntos porcentuales el tipo marginal máximo de la base liquidable del ahorro, que pasa del $28\%$ al $30\%$ para la parte que exceda de $300.000€$, con efectos desde el 1 de enero de 2025. Esto último reduce (sin eliminar) la brecha entre el tipo marginal máximo de la renta general (que en algunas CCAA supera el $47-50\%$) y el de la renta del ahorro, lo que sigue siendo un argumento vivo en el debate sobre la equidad horizontal del modelo dual.

\subsection{Marco jurídico actual: Ley 35/2006}
Conocido el desarrollo histórico del modelo, veamos como se configura en la actualidad. 

Como decíamos, el IRPF vigente se regula con carácter principal en la ya mencionada LIRPF 35/2006, desarrollada reglamentariamente por el Real Decreto 439/2007, de 30 de marzo (Reglamento del IRPF). Como en todo tributo, rige supletoriamente la Ley 58/2003, General Tributaria, para cuanto no esté específicamente regulado en la normativa del impuesto (procedimientos de gestión e inspección, infracciones y sanciones, revisión en vía administrativa). 

A esto se añade un tercer nivel normativo, específico del IRPF por su condición de tributo cedido parcialmente: la Ley 22/2009, de 18 de diciembre, que regula el sistema de financiación de las Comunidades Autónomas y determina el alcance de su capacidad normativa sobre el impuesto. Aunque esto último será tratado más adelante, que conviene tenerlo presente ya que la LIRPF no agota el marco jurídico aplicable: hay que sumarle, para cada contribuyente, la normativa de su CA de residencia.

\subsubsection{Ámbito de aplicación: ¿Dónde se aplica el impuesto?}
El IRPF se aplica en todo el territorio español, sin perjuicio de los regímenes forales de Concierto Económico (País Vasco) y Convenio Económico (Navarra).\footnote{Las características particulares de Canarias, Ceuta y Melilla en otros tributos no tienen cabida en el IRPF, ya que no son hay cambios sustanciales de base imponible, a diferencia del IVA/IGIC.
}

El elemento verdaderamente decisivo del ámbito de aplicación, sin embargo, no es tanto el territorial en sentido físico como la determinación de quién queda sujeto por obligación personal (renta mundial) frente a quién queda sujeto solo por obligación real (renta de fuente española, vía IRNR). Esa frontera la traza el concepto de \textbf{residencia habitual} del art. $9$ LIRPF. Se entiende que el contribuyente tiene su residencia habitual en territorio español cuando concurra cualquiera de estas circunstancias, aplicadas de forma independiente:
\begin{la}
    \item \textbf{Marco temporal}. Que permanezca más de $183$ días durante el año natural en territorio español, computándose las ausencias esporádicas salvo que se acredite residencia fiscal en otro país (y, si el destino es un paraíso fiscal, la Administración puede exigir la prueba de permanencia efectiva); o,
    \item \textbf{Intereses económicos}. Que radique en España el núcleo principal o la base de sus actividades o intereses económicos, de forma directa o indirecta. 
\end{la}

Aunque la ley los formula como alternativos, la doctrina y cierta jurisprudencia menor han entendido que operan con un cierto orden de prelación práctico: primero se analiza la permanencia física, y solo si no resulta concluyente se acude al centro de intereses económicos. 

Además, a estos dos criterios positivos se suma una presunción iuris tantum de residencia: si no se acredita lo contrario, se presume que el contribuyente reside en España cuando residan aquí, de forma habitual, su \textbf{cónyuge} no separado legalmente y los \textbf{hijos menores} de edad que dependan de él; presunción que, por su naturaleza, puede desvirtuarse con prueba en contrario cuya carga se traslada al contribuyente. Como es de suopner, este último criterio es el que más litigiosidad genera en la práctica. Por ejemplo en los que se conoce como residentes fiscales de fin de semana, ya que combinan un elemento objetivo (dónde vive la familia) con una probática compleja (viajes, contratos de arrendamiento, consumos, tarjetas de crédito) que la Inspección utiliza habitualmente para reconstruir calendarios de presencia.

Dentro del ámbito de aplicación conviene situar también dos figuras que iluminan bien la lógica del legislador. En primer lugar, el \textbf{régimen especial de trabajadores desplazados} (art. $93$ LIRPF, popularmente Ley Beckham), que permite a quien traslada su residencia fiscal a España optar por tributar conforme a las reglas del IRNR (tipo fijo, escala no progresiva) manteniendo formalmente la condición de contribuyente del IRPF, durante el periodo impositivo del cambio y los cinco siguientes. Además, tras la reforma operada por la Ley 28/2022 (Ley Startups), este ámbito subjetivo se amplió sustancialmente ya que dejó de limitarse al trabajador desplazado, y se extendió también a teletrabajadores internacionales, administradores de determinadas entidades y emprendedores, quedando eso sí excluidos los deportistas profesionales desde 2015.\footnote{Como vemos, el ámbito de aplicación no es una cuestión meramente técnica sino un instrumento consciente de política fiscal (competencia entre jurisdicciones por atraer capital humano cualificado) y, a menudo, es un sustitutivo elegante a los gastos fiscales directos, ya que formalmente no es un gasto fiscal en sentido estricto sino un régimen especial de sujeción.
} En segundo lugar, detsca también el \textbf{régimen de cambio de residencia a paraísos fiscales} (art. $8.2$ LIRPF): un español que traslade su residencia a un territorio calificado como paraíso fiscal no pierde la condición de contribuyente por el IRPF durante el periodo del cambio y los cuatro siguientes, creándose de esta forma una \textit{cláusula antiabuso} frente a la deslocalización fiscal oportunista.

\subsubsection{Sujeto pasivo (contribuyente)}
Técnicamente, la LIRPF ya no habla de \textit{sujeto pasivo} sino de contribuyente (art. $8$ LGT y art. $8$ LIRPF), pues el IRPF, al ser un impuesto directo sin repercusión, hace coincidir en la misma persona al sujeto pasivo y al que soporta la carga económica del tributo, distinción que la diferencia conceptualmente frente a las figuras indirectas.

Atendiendo a la LIRPF, son contribuyentes:
\begin{la}
    \item Las personas físicas con residencia habitual en territorio español conforme al art. $9$ ya visto; y,
    \item Por extensión, las personas físicas de nacionalidad española con residencia habitual en el extranjero en determinados supuestos tasados del art. $10$ LIRPF. Esencialmente, miembros de misiones diplomáticas y funcionarios públicos españoles destinados en el extranjero, salvo que ya tuvieran esa condición de no residentes con anterioridad a su destino.
\end{la}

El lector atento habrá notado que la unidad familiar (art. $82-84$ LIRPF) no se menciona como sujeto pasivo del impuesto. Esto es consecuencia directa de que, por lo que a la LIRPF respecta, no tiene personalidad jurídico-tributaria propia, sino que se configura como una opción de tributación: cada miembro sigue siendo individualmente contribuyente, pero pueden optar, año a año, por acumular sus rentas en una única declaración con reglas específicas.\footnote{Entre las más descadas, la aplicación de una escala conjunta sobre la renta acumulada, con una reducción en la base imponible general de $3.400$ euros para matrimonios y $2.150$ euros para familias monoparentales, como corrección parcial del efecto de la progresividad sobre la renta acumulada.
} Como sabemos, es el legado directo de la STC 45/1989: la ley no puede imponer la acumulación, solo ofrecerla como opción y corregir su efecto agravante.

Por último, el art. $8.3$ LIRPF excluye expresamente del concepto de contribuyente a las sociedades civiles no sujetas al Impuesto sobre Sociedades, herencias yacentes, comunidades de bienes y demás entidades del art. $35.4$ LGT: estas entidades no tributan como tales, sino que sus rentas se atribuyen a los socios, herederos, comuneros o partícipes conforme al régimen de atribución de rentas (arts. $8.3$ y $86-90$ LIRPF), que es la técnica que evita la doble imposición de una renta generada colectivamente pero carente de sujeto tributario propio.

\subsubsection{Hecho imponible, exenciones y periodo impositivo}
El art. $6.1$ LIRPF define el hecho imponible como la obtención de renta por el contribuyente; una definición deliberadamente amplia y residual, coherente con la vocación sintética del impuesto. No obstante, el propio art. $6.2$ LIRPF concreta los componentes de esa renta listando las cinco fuentes que el impuesto reconoce: rendimientos del trabajo, rendimientos del capital (mobiliario e inmobiliario), rendimientos de actividades económicas, ganancias y pérdidas patrimoniales, e imputaciones de renta que la propia ley establezca (por ejemplo, inmuebles no arrendados o el régimen de transparencia fiscal internacional). Esta clasificación por fuentes es la que luego, en el esquema de liquidación, determinará si una renta se integra en la base general o en la base del ahorro; y es también uno de los principales puntos de tensión con el ideal teórico, que no distingue por fuente.

Un matiz técnico que conviene tener muy asentado es que el hecho imponible se predica sobre la renta mundial del contribuyente residente, a diferencia del IRNR, que grava solo la renta de fuente española obtenida por no residentes bajo el criterio de territorialidad.\footnote{Tal y como establece el art. $2$ LIRPF: \textit{constituye el objeto de este impuesto la renta del contribuyente, entendida como la totalidad de sus rendimientos, ganancias y pérdidas patrimoniales y las imputaciones de renta que se establezcan por la ley, con independencia del lugar donde se hubiesen producido y cualquiera que sea la residencia del pagador})
} Esta es la contrapartida lógica de la sujeción por residencia: quien tributa por obligación personal lo hace por su capacidad económica global, no solo por la generada en territorio nacional.

En cuanto a las exenciones (art. $7$ LIRPF), es una lista extensa con más de treinta supuestos que no conviene memorizar exhaustivamente sino agrupar por finalidad:
\begin{la}
    \item \textbf{Capacidad económica real}. Exenciones que responden a una lógica de no capacidad económica real o de compensación de un daño (estructurales en sentido amplio). Por ejemplo, indemnizaciones por despido dentro de los límites del Estatuto de los Trabajadores, indemnizaciones por responsabilidad civil por daños personales, prestaciones por incapacidad permanente absoluta o gran invalidez de la Seguridad Social.
    \item \textbf{Objetivo redistributivo}. Exenciones con finalidad social o de política de rentas (más próximas al concepto de gasto fiscal que desarrollaremos en el apartado 3). Por ejemplo, exención de becas públicas, anualidades por alimentos percibidas de los padres por decisión judicial, prestaciones por maternidad/paternidad, determinadas ayudas a víctimas de violencia de género o de terrorismo.
    \item \textbf{Incentivos al ahorro}. Exenciones técnicas, dirigidas a evitar dobles imposiciones o a facilitar la movilidad del ahorro. Por ejemplo, dividendos y participaciones en beneficios exentos hasta ciertos límites históricamente o la exención por reinversión en vivienda habitual para mayores de $65$ años.
\end{la}

Un dato curioso en este sentido es que el Salario Mínimo Interprofesional (SMI), contrario a lo que se cree, no está formalmente exento en el art. $7$ LIRPF, sino que queda libre de tributación efectiva gracias al juego combinado del mínimo personal y familiar y de la reducción por obtención de rendimientos del trabajo. Es decir, técnicamente tributa, pero su cuota resultante es cero. Es un matiz útil que debe impedir que nos precipitemos al error, muy extendido incluso entre profesionales, de decir que el SMI está exento.

Por lo que respecta al periodo impositivo, el art. $12$ LIRPF fija la regla general en el \textbf{año natural}, con devengo el 31 de diciembre de cada año; coincidiendo, por tanto, hecho imponible y periodo impositivo en un impuesto periódico de devengo instantáneo al final del periodo. La única excepción relevante es el fallecimiento del contribuyente en un día distinto (art. $13$ LIRPF), que determina un periodo impositivo inferior al año natural. En estos casos, se devengará el impuesto en la fecha del fallecimiento y, por tanto, se mantiene la coherencia con el principio de personalidad del impuesto: la obligación tributaria no puede sobrevivir, en su hecho imponible, a la persona física que la genera. El \textbf{elemento temporal} del impuesto también suscita mucha confusión en lo relativo al cómputo de las rentas, en particular la compleja normativa de imputación temporal del art. $14$ LIRPF, que es donde se libran buena parte de las controversias prácticas entre Administración y contribuyente.\footnote{Principalmente, el criterio de devengo como regla general, con reglas especiales para rentas no satisfechas por resolución judicial, rentas obtenidas en el extranjero con restricciones a la transferencia de divisas, o ganancias patrimoniales derivadas de ayudas públicas, etc.
}

\subsubsection{Gestión y aplicación del impuesto}
La gestión y aplicación del IRPF corresponde, con carácter general, a la Agencia Estatal de Administración Tributaria (AEAT), salvo en los territorios forales, donde corresponde a las respectivas Haciendas Forales. El sistema se articula sobre tres pilares:
\begin{lnum}
    \item \textbf{Pilar 1: Autoliquidación}. La autoliquidación, prevista en el art. $97$ LIRPF, diferencia al IRPF de otros tributos que cuentan con liquidación administrativa. En el IRPF es el propio contribuyente quien calcula su deuda tributaria y la ingresa, mediante el Modelo 100, en el plazo que anualmente fija una Orden Ministerial (habitualmente entre abril y finales de junio del ejercicio siguiente);
    \item \textbf{Pilar 2: Pagos a cuenta}. Los pagos a cuenta, regulados en los arts. $99-101$ LIRPF, son retenciones sobre rendimientos del trabajo y del capital, ingresos a cuenta y pagos fraccionados de actividades económicas, que operan como anticipos de la cuota final y cuya función, más allá de recaudatoria, es también de control (permiten a la Administración cruzar información con terceros pagadores) y de suavización de la carga financiera para el contribuyente, al fraccionar en el tiempo un pago que de otro modo se concentraría en un único momento; y,
    \item \textbf{Pilar 2: Propuesta de declaración}. El borrador de declaración (art. $98$ LIRPF), confeccionado por la Administración a partir de la información de que dispone (retenciones, datos fiscales comunicados por terceros) y que el contribuyente puede confirmar, modificar o rechazar. Merece mención porque ilustra un cambio de paradigma en la gestión de masas. Desde un modelo en que el contribuyente calcula y declara, se ha ido evolucionando hacia un modelo de declaración asistida en el que la Administración anticipa la mayor parte de la información con los datos que ya obran en su poder (retenciones, movimientos inmobiliarios, datos bancarios comunicados vía suministro de información), reduciendo el margen de error material del contribuyente pero también, según ha señalado la doctrina, generando un cierto riesgo de confirmación pasiva de datos incorrectos.
\end{lnum}

Debemos tener en cuenta que no todos los contribuyentes están obligados a declarar. Quedan excluidos, con carácter general, quienes obtengan rendimientos íntegros del trabajo inferiores a $22.000$ euros anuales de un único pagador (o $15.876$ euros si proceden de más de un pagador y la suma del segundo y restantes supera $1.500$ euros), así como quienes obtengan rendimientos del capital mobiliario y ganancias patrimoniales sometidos a retención por un importe conjunto inferior a $1.600$ euros anuales. Esta arquitectura de umbrales no es casual, ya que revela que el legislador combina un impuesto formalmente universal (todo residente es contribuyente) con una gestión de masas que exime de declarar a quien ya ha satisfecho su obligación material vía retención, optimizando así el coste de cumplimiento sin sacrificar recaudación.

\subsubsection{Distribución competencial: Estado, Comunidades Autónomas y territorios forales}
Aquí es donde conviene aplicar con rigor la distinción que hacíamos al principio. El IRPF es, formalmente, un tributo estatal cedido parcialmente a las CCAA de régimen común, al amparo de la LOFCA (8/1980) y desarrollado por la Ley 22/2009. Esta cesión opera en dos planos distintos, que no deben confundirse:
\begin{la}
    \item \textbf{Cesión del rendimiento}. El $50\%$ de la recaudación líquida del IRPF corresponde a la CA de residencia habitual del contribuyente (art. $72$ LIRPF fija los criterios de residencia autonómica, con reglas de desempate cuando la permanencia no es concluyente, entre ellas la de mayor base imponible);
    \item \textbf{Cesión de capacidad normativa parcial}. Las CCAA pueden regular, dentro de los límites fijados por la Ley 22/2009, la escala autonómica aplicable a la base liquidable general (que se suma a la escala estatal), el importe del mínimo personal y familiar (dentro de un margen del $\pm10\%$ sobre las cuantías estatales), y determinadas deducciones propias en la cuota. No pueden, en cambio, regular la base imponible, el hecho imponible, la escala del ahorro ni los aspectos de gestión del impuesto, de ahí que sea técnicamente incorrecto hablar de gestión compartida.
\end{la}

Esta arquitectura explica un fenómeno muy debate en las esferas políticas nacionales: la carga fiscal del IRPF varía sensiblemente entre CCAA, pese a tratarse formalmente del mismo impuesto. Como vemos, la diferencia real no está en el hecho imponible ni en la base, que son idénticos en todo el territorio de régimen común, sino en la escala autonómica y en las deducciones propias, ejercicio de la capacidad normativa cedida.

Distinto es el régimen de los territorios forales (País Vasco, por el Concierto Económico aprobado por Ley 12/2002, y Navarra, por el Convenio Económico aprobado por Ley 28/1990). En estas CCAA no hay cesión parcial, sino una exacción, gestión, liquidación, recaudación e inspección íntegramente forales del IRPF, con capacidad normativa propia sujeta a los principios de armonización del Concierto/Convenio, con estipulaciones como que la presión fiscal efectiva global no sea inferior a la del territorio común, respeto a la libertad de circulación, etc. Por tanto, es un régimen de naturaleza distinta al de cesión, y conviene marcarlo con claridad porque es un error frecuente calificarlas como de \textit{régimen especial}, cuando en realidad es un sistema de concierto bilateral de naturaleza histórica cualitativamente distinto de la financiación autonómica común.

\begin{specialblock}{\textbf{País Vasco:} No una Comunidad Autonómica más, sino tres}
    Ya fijamos la diferencia cualitativa entre la cesión parcial del régimen común y el Concierto Económico del País Vasco y el Convenio de Navarra. Lo que conviene añadir aquí, porque afecta directamente al catálogo de gastos fiscales y no solo a la arquitectura competencial, es un matiz que sorprende a quien no conoce el sistema en detalle: el País Vasco no tiene un único IRPF foral, sino tres. 

    La competencia tributaria corresponde, dentro de Euskadi, a las Instituciones de cada uno de sus tres Territorios Históricos —Álava/Araba, Bizkaia y Gipuzkoa—, cada uno con su propia Diputación Foral, su propia Norma Foral del IRPF, y su propia tarifa, mínimos y deducciones, coordinadas entre sí a través del Órgano de Coordinación Tributaria de Euskadi (creado por la Ley 3/1989 del Parlamento Vasco) y sometidas a un alto grado de armonización voluntaria desde la reforma conjunta de 2014, pero sin identidad normativa plena entre ellas. [Probable] Navarra, por su parte, mantiene un único IRPF foral para todo su territorio, regulado por Ley Foral propia al amparo del Convenio Económico. La consecuencia práctica es que, en sentido estricto, no existen "dieciocho versiones del IRPF" en España (diecisiete Comunidades Autónomas más el Estado), sino potencialmente más: territorio común con sus diecisiete variantes autonómicas, más Navarra, más tres regímenes forales vascos coordinados pero no idénticos entre sí.
\end{specialblock}

\subsection{Esquema básico de liquidación}
Cerraremos este bloque con el esquema que sirve de columna vertebral para el resto del tema. Lo presentamos en su secuencia lógica, deteniéndonos en los puntos donde se sitúan las divergencias que sistematizaremos más adelante:
\begin{lnum}
    \item \textbf{Rendimientos íntegros}. Por cada una de las fuentes del art. $6.2$ LIRPF (trabajo, capital mobiliario, capital inmobiliario, actividades económicas), más las imputaciones de renta y las ganancias y pérdidas patrimoniales;
    \item \textbf{Rendimiento neto}. Sobre el punto anterior, se aplican los gastos deducibles propios de cada fuente (cotizaciones a la Seguridad Social en el trabajo, gastos de conservación y financiación en el capital inmobiliario, gastos deducibles conforme a la normativa del IS en actividades económicas) y, en su caso, las reducciones específicas (reducción por obtención de rendimientos del trabajo, reducción por rendimientos irregulares o generados en más de dos años, con el $30\%$ aplicable con límite de $300.000$ euros anuales);
    \item \textbf{Integración y compensación}. En este punto se materializa la dualización nacida en 2006: (i) los rendimientos del trabajo, capital inmobiliario y actividades económicas, junto con determinadas imputaciones, se integran en la base general; (ii) los rendimientos del capital mobiliario (salvo excepciones) y las ganancias y pérdidas patrimoniales derivadas de transmisiones se integran en la base del ahorro. Cada una con reglas propias de compensación entre rendimientos positivos y negativos y con límites cruzados entre ambas bases;
    \item \textbf{Base imponible: dual}. Se diferencian la base imponible general y la base imponible del ahorro, resultado de las operaciones anteriores;
    \item \textbf{Reducciones en la base imponible}. Fundamentalmente, aportaciones a sistemas de previsión social (planes de pensiones, con límite conjunto de $1.500$ euros anuales, ampliable hasta $8.500$ euros si proceden de contribuciones empresariales), pensiones compensatorias y anualidades por alimentos al cónyuge;
    \item \textbf{Base liquidable: dual}. Se determina la base liquidable general y la base liquidable del ahorro, diferenciándolas;
    \item \textbf{Mínimo personal y familiar} (arts. $56-61$ LIRPF). Se aplican los mínimos personal y familiar a las bases imponibles. $5.550€$ anuales con carácter general, incrementado por edad (mayores de $65$ y $75$ años), por descendientes, ascendientes y por discapacidad. Este es, técnicamente, el mecanismo mediante el cual el legislador excluye de gravamen la renta que se estima necesaria para atender las necesidades vitales básicas del contribuyente y su familia. Como veremos más adelante, es la primera conexión con HAIG-SIMONS: un mínimo exento fundado en necesidades vitales no es, en puridad, un gasto fiscal en el sentido de renuncia recaudatoria con fin extrafiscal, sino un elemento estructural del propio concepto de capacidad económica gravable;
    \item \textbf{Cuota íntegra estatal y autonómica}. Aplicación de la escala progresiva de la base general (suma de escala estatal $+$ escala autonómica) y de la escala de la base del ahorro (también progresiva por tramos desde 2010, aunque proporcional en su origen de 2006);
    \item \textbf{Cuota líquida}. Aplicación de las deducciones en cuota (vivienda habitual en régimen transitorio, donativos, inversión en empresas de nueva creación, deducciones autonómicas, maternidad, familia numerosa); y,
    \item \textbf{Cuota diferencial}. Cuota líquida minorada en las retenciones y pagos a cuenta ya soportados durante el ejercicio, dando lugar al resultado a ingresar o a devolver.
\end{lnum}

Gráficamente, aunque no es formalmente el esquema de liquidación habitual, vamos a apoyarnos en el siguiente:

\imagenfit{Fig1.png}{Esquema de liquidación: Impuesto sobre la Renta de las Personas Físicas}{Elaboración propia}
















\clearpage
\section{Divergencias estructurales: ¿Cómo se define la renta gravable en España?}
\puntosclave{
    
}


Antes de que existiera un concepto unificado de renta a efectos fiscales, la doctrina hacendística clásica manejaba dos formas muy distintas de entender qué cuenta como renta gravable:
\begin{la}
    \item \textbf{Teoría de la fuente} (renta-producto). Solo es renta el flujo periódico y regular que procede de una fuente productiva estable y duradera —un salario, el alquiler de un inmueble, el beneficio recurrente de un negocio—. Bajo esta óptica, una plusvalía esporádica por vender un piso, un premio de lotería, o una herencia recibida no serían renta, precisamente por no ser el "fruto" periódico de una fuente productiva, sino un hecho puntual y no recurrente; y,
    \item \textbf{Teoría del incremento patrimonial neto} (renta-entrada, accretion theory). Frente a esa visión restrictiva, esta segunda teoría considera renta cualquier incremento del poder económico del sujeto, tenga o no carácter periódico y proceda de donde proceda —trabajo, capital, una donación recibida, un premio, una plusvalía—. Lo relevante no es la fuente ni la regularidad, sino el simple hecho de que la capacidad económica del sujeto ha aumentado.
\end{la}

A lo largo de los años, diversos economistas se han centrado en la cuestión de cuá sería el ideal teórico para establecer como referencia en el diseño de una figura impositiva sobre la renta. La segunda corriente, no obstante, es la que a principios del siglo pasado surgió como vencedora: la renta SHS es la culminación técnica de esta segunda tradición. Fruto del trabajo de tres economistas, que la formularon en momentos distintos: SCHANZ (1896), HAIG (1921) y SIMONS (1938), su definición se resume en la fórmula:

$$Y=C+\Delta W$$

Es decir: renta es igual a consumo m´s variación neta del patrimonio del sujeto. De esa fórmula se derivan tres exigencias que un impuesto sobre la renta ideal debería cumplir:
\begin{la}
    \item \textbf{Comprehensividad de fuente}. Da igual de dónde venga el dinero, todo tributa igual (salario, dividendo, herencia, premio);
    \item \textbf{Devengo, no realización}. Se grava el incremento de valor cuando se produce, no cuando se materializa en una venta (lo que implica gravar también las plusvalías latentes de activos que aún no se han vendido);
    \item \textbf{Inclusión de rentas en especie e imputadas}. El valor de uso de bienes propios (como vivir en tu propia vivienda sin pagar alquiler) también es renta.
\end{la}

Este es el aparato conceptual previo necesario para abordar nuestra pregunta, puesto que ahora no nos preguntaremos \textit{qué es la renta en términos económicos} ([\authorfont{Ver Tema 4.B.12}]), sino \textbf{cómo diverge, de forma sistemática, de ese ideal teórico}. Pregunta tiene una respuesta metodológica concreta, y es con la que empezaremos.

\subsection{Benchmarking fiscal: De la teoría al método}
La respuesta la dió SURREY (1967), entonces Secretario Adjunto del Tesoro de Estados Unidos para Política Fiscal, en un artículo seminal que desarrolló después en su libro (\textit{Pathways to Tax Reform}, 1973): para poder afirmar que una norma tributaria constituye una desviación respecto de un impuesto de renta ideal, hace falta fijar previamente cuál es el sistema tributario de referencia o benchmark frente al que se mide esa desviación. 

SURREY llamó gastos fiscales a las disposiciones que se apartan del benchmark con una finalidad de fomento económico o social equiparable a la de un programa de gasto público directo, y cuya diferencia con este último es, en última instancia, solo instrumental: en vez de dar una subvención, el Estado renuncia a una parte de la recaudación que le correspondería conforme al benchmark. 

Conviene no ocultar, que la elección del benchmark no es una operación neutral ni técnicamente unívoca: el propio SURREY reconocía que un benchmark demasiado estricto (SHS puro) califica como gasto fiscal casi cualquier rasgo estructural del diseño real del impuesto (incluida la propia progresividad de la tarifa), mientras que un benchmark demasiado laxo no deja ninguna desviación que analizar.\footnote{
Por ejemplo, tomar el propio diseño legal vigente como referencia de sí mismo. La literatura posterior a SURREY, como MCDANIEL y SURREY (1985), y en el ámbito de la OCDE los sucesivos informes sobre gasto fiscal, ha convivido con esta tensión sin resolverla del todo, optando en la práctica por benchmarks intermedios: normativos, pero no puramente SHS.
} 

\subsubsection{Método en España: Presupuesto de Beneficios Fiscales}
España no ha sido ajena a esta cuestión metodologíca, aunque la aplica con un benchmark propio y no siempre explícito. El fundamento constitucional está en el art. 134.2, que dispone que los PGE consignarán el importe de los beneficios fiscales que afecten a los tributos del Estado.

El Presupuesto de Beneficios Fiscales (PBF) se elabora en España desde 1979, aunque la obligación de acompañarlo de una Memoria explicativa de su cuantificación no se formalizó hasta Ley 41/1994 (PGE-95). Desde ese año, cada año el proyecto de Ley de Presupuestos incorpora esta Memoria que cuantifica, tributo por tributo, el coste recaudatorio de cada beneficio fiscal.

Ahora bien, conviene ser ligeramente crítico con el método español. La Memoria no parte de un benchmark SHS puro, sino de un sistema tributario de referencia que ya da por asumidas ciertas características estructurales del IRPF vigente y que solo trata como beneficio fiscal lo que se aparta de ese diseño ya asumido como normal. Por tanto, es imperativo en primer lugar conocer esas divergencias en la sombra que no computan como gasto fiscal.

Distinguimos, en especial, dos: temporal y de fuente. Pasemos a analizarlas.\footnote{Esta distinción, en conjunto, es la que vertebra justifica la división que hemos acometido: divergencias estructurales y gastos fiscales. Fijado el benchmark (SHS) y el método, analicemos cuatro grandes divergencias estructurales respecto de la renta SHS, organizadas no por la naturaleza económica de la divergencia:

\textcolor{red}{1. **Divergencia temporal**: realización frente a devengo en la tributación de las ganancias patrimoniales.
2. **Divergencia por fuente**: la estructura dual (renta general / renta del ahorro), con tratamiento conjunto, dentro de este bloque, de la imputación de rentas y las rentas en especie, en la medida en que —como tú mismo apuntas— son, en puridad, un caso particular de renta no percibida en efectivo que el legislador decide gravar o no gravar según su origen.
3. **Debates abiertos sobre la arquitectura del impuesto**: de un lado, la exclusión de las transmisiones gratuitas (herencias, donaciones) del hecho imponible del IRPF por su derivación al ISD, y la pregunta más amplia sobre qué otras figuras podrían integrarse en un impuesto sobre la renta más comprehensivo y por qué razones —sociales, culturales, de viabilidad política— no lo están; de otro, el debate renta-consumo y la neutralidad intertemporal del ahorro, con particular atención al tratamiento fiscal de la previsión social.
4. **Balance final**: caracterización del IRPF español dentro de la tipología comparada de impuestos sobre la renta (sintético puro / dual / híbrido), como cierre valorativo de todo el apartado.}
}

\subsection{Divergencia temporal: Realización frente a devengo}
Como hemos visto, la definición SHS exige gravar la renta cuando se genera, con independencia de que se materialice en efectivo. Si el valor de un activo aumenta, ese incremento es renta en el periodo en que se produce. 

El IRPF español, como prácticamente todos los sistemas comparados, no sigue este criterio para las ganancias y pérdidas patrimoniales: el art. $33.1$ LIRPF define la ganancia o pérdida patrimonial como la variación en el valor del patrimonio del contribuyente que se ponga de manifiesto con ocasión de cualquier alteración en su composición, salvo que la propia ley disponga lo contrario. Es decir, por lo general exige un negocio jurídico que materialice el cambio (venta, permuta, donación, expropiación. Dicho de otro modo: mientras el contribuyente conserve el activo, por mucho que se revalorice, no hay hecho imponible. 

El impuesto español adopta, pues, el \textbf{criterio de realización}.

\subsubsection{Argumentos a favor: ¿razones de la divergencia?}
Como podemos suponer, esta no es una divergencia arbitraria. Por el contrario, responde a razones defendibles aunque no exentas de coste. Así que, ¿por qué se ecoge este criterio en contra del ideal?

En primer lugar, se sostiene bajo el argumento de la \textbf{iliquidez}. Gravar por devengo obligaría al contribuyente a pagar impuesto, en efectivo, sobre una renta que no ha percibido de la misma forma, lo que puede forzar ventas no deseadas solo para hacer frente al pago (un problema especialmente agudo en activos ilíquidos: inmuebles, participaciones en empresas familiares no cotizadas). Por ejemplo, imaginemos el propietario de acciones que se revalorizan, éste tendría que encontrar liquidez ajena al propio activo para satisfacer la deuda tributaria. Por otro lado, el devengo exige \textbf{valorar anualmente} todo el patrimonio del contribuyente a precio de mercado, lo cual es sencillo para activos cotizados pero extraordinariamente costoso, y en muchos casos arbitrario, para activos sin mercado organizado (inmuebles, empresas no cotizadas, obras de arte, etc.). Por último, suele defenderse por razones de \textbf{simplicidad administrativa y cumplimiento}. La realización ata el hecho imponible a un acontecimiento jurídico verificable (una escritura, un contrato), lo que facilita tanto la autoliquidación por el contribuyente como el control por la Administración, frente a la incertidumbre inherente a una valoración anual estimada.

Aunque estas razones son sólidas para activos ilíquidos y de valoración incierta, su sustento es mucho peor cuando se trata de activos financieros líquidos y cotizados diariamente, donde la valoración diaria existe y es pública. La literatura más actual defiende que la justificación última y real para establecer un criterio de realización es de coste administrativo y restricción de liquidez, no un argumento de principio sobre qué es renta. Por ello, muchos autores defienden en documentos o informes de reforma fiscal, soluciones diferenciadas por tipo de activo en vez de una regla uniforme (como la española).

Además, se han dado algunos argumentos económicos en contra de la utilización de este criterio.

\paragraph{Ineficiencia: Lock-in effect}
En primer lugar, al condicionar el gravamen a la venta, el sistema crea un incentivo a no vender activos revalorizados y diferir indefinidamente el pago del impuesto. Como podemos suponer, esto distorsiona las decisiones de inversión: un contribuyente puede mantener una cartera subóptima solo para evitar materializar la plusvalía, reduciendo la movilidad del capital hacia sus usos más productivos. 

Constituye ésta uno de los costes de ineficiencia más citados en la literatura y explica por qué muchas propuestas de reforma (mark-to-market taxation) se dirigen precisamente a activos líquidos, donde el coste de este efecto es mayor y su corrección más viable técnicamente.

\paragraph{Inflación: Corrección monetaria} 
Hasta 2014, la LIRPF contemplaba dos mecanismos que corregían el desajuste entre ganancia nominal y ganancia real, cada uno con sus objetivos y su método particular: (i) los \textbf{coeficientes de corrección monetaria}, que actualizaban el valor de adquisición de los inmuebles conforme a la depreciación del poder adquisitivo de la moneda; y, (ii) \textbf{los coeficientes de abatimiento}, que reducían la ganancia gravable en función de la antigüedad de los bienes comprados antes del 31/12/1994 (como transición al endurecimiento del régimen de plusvalías). 

No obstante, la Ley 26/2014 suprimió íntegramente los coeficientes de corrección monetaria y limitó drásticamente los coeficientes de abatimiento. Como podemos intuir, el segundo no es especialmente importante, por estar configurado para un contexto histórico particular. Sin embargo, el primero sí que es interesante en términos económicos. 

Imagina que compraste un piso en el año 2000 por $100.000€$, y lo vendes ahora, en 2024, por $200.000€$. Estos $100.000€$ son lo Hacienda llama ganancia nominal: la simple resta entre lo que pagaste y lo que has cobrado. Pero, como economistas, sabemos que eso no es así. Esos $100.000€$ del año 2000 no valían lo mismo que $100.000€$ de hoy. Con la inflación acumulada, para comprar hoy lo mismo necesitarías bastante más. Pongamos, solo como ejemplo ilustrativo, que ese importe correspone hoy a unos $170.000€$. Es decir: si tu piso vale ahora $200.000€$, y solo mantener el poder adquisitivo de tu dinero original ya costaría $170.000€$, tu ganancia real son solo $30.000€$. El resto, es inflación, no renta.

Por tanto, antes de 2015, la ley tenía un mecanismo que hacía este ajuste: actualizaba el precio de compra según la inflación acumulada antes de calcular la ganancia. Esta supresión, motivada por razones recaudatorias en el contexto de consolidación fiscal post-crisis, tiene una implicación conceptual importante: desde 2015, el IRPF grava la ganancia nominal, no la ganancia real, lo que en periodos de inflación elevada supone gravar como renta una parte de lo que en realidad es solo depreciación monetaria del capital invertido. Otra divergencia adicional respecto de un benchmark entendido en términos de poder adquisitivo real (SHS), y no meramente nominal.

\paragraph{Neutralidad intertemporal débil: Asimetría en el tratamiento de pérdidas}
Por otro lado, la lógica de neutralidad de SHS exigiría un tratamiento simétrico de ganancias y pérdidas: si se grava íntegramente el incremento, debería compensarse íntegramente la pérdida. 

No obstante, el IRPF limita esa simetría: el saldo negativo de rendimientos del capital mobiliario solo puede compensarse con el saldo positivo de ganancias y pérdidas patrimoniales de esa misma base (y viceversa) con un límite porcentual sobre el saldo positivo, límite que se ha ido ampliando progresivamente desde su introducción en la reforma de 2015 hasta el $25\%$ actual. 

Pensémos en sus implicaciones. Imaginemos que en un mismo año ganamos $10.000€$ vendiendo acciones de la empresa A (ganancia patrimonial) y perdemos $8.000€$ con un depósito o unos bonos que rindieron menos de lo esperado (rendimiento negativo del capital mobiliario).

Si el sistema fuera plenamente simétrico, esos $8.000€$ de pérdida deberían poder restarse sin más de los $10.000€$ de ganancia, dejándote con solo $2.000€$ netos por los que tributar. Pero la ley pone un tope: esa pérdida solo puede compensar hasta un $25\%$ del saldo positivo de ese mismo cajón (la base del ahorro) en ese año. Si el $25\%$ de tus ganancias del año no es suficiente para absorber toda tu pérdida, el resto no se pierde, pero tienes que esperar a compensarlo en años futuros, poquito a poco.

¿Por qué importa esto? Porque significa que si tus inversiones son volátiles acabas pagando más impuesto en total, a lo largo de los años, que alguien que tuvo exactamente la misma rentabilidad media pero de forma estable año tras año. El fisco te grava el $100\%$ de lo que ganas cada año, pero solo te deja descontar una parte de lo que pierdes cada año. No es una regla neutral: penaliza sutilmente el riesgo y la irregularidad, y además te obliga a "prestarle" a Hacienda, sin interés, la parte de la pérdida que todavía no has podido compensar.

Es una asimetría adicional, de menor calado que las anteriores pero coherente con el mismo patrón: cuando la aplicación estricta de la neutralidad SHS entra en conflicto con la protección de la recaudación, el legislador español ha resuelto sistemáticamente a favor de la segunda. 

\subsection{Divergencia de fuente: Dualización del impuesto}
Otra de las divergencias claras, y más visibles, respecto al benchmark teórico es la divergencia de fuente de rentas. Sabemos que SHS exige que la renta tribute igual con independencia de su origen: un euro de salario, un euro de dividendo y un euro de plusvalía son, en términos de capacidad económica, exactamente lo mismo, y deben tributar al mismo tipo. 

El IRPF español vemos como se aparta de SHS en, al menos, tres planos acumulativos: (i) una brecha estructural de tipos entre base general y base del ahorro, con fundamento en la movilidad del capital pero con coste en equidad horizontal y vertical; (ii) un incentivo persistente al arbitraje de calificación; y, (iii) una asimetría adicional en el tratamiento de las rentas en especie e imputadas, que depende más de la forma en que se percibe la renta que de su sustancia económica. Esta es, en esencia, la \textbf{ruptura del principio de comprehensividad de fuente}. 

\subsubsection{Debate: Modelo dual y fundamentos}
Estas tres manifestaciones no son incidentes aislados de la norma, sino consecuencia directa y predecible de haber optado por el modelo dual en lugar del sintético puro.

Lo primero que debemos saber es que España no inventó este diseño. Lo importó, con variaciones, del Dual Income Tax (DIT) desarrollado por los países nórdicos a partir de finales de los años ochenta (Dinamarca, 1987; Noruega, 1992; Suecia, 1991), sistematizado después en la literatura académica por autores como SØRENSEN y CNOSSEN. El argumento a favor del modelo dual descansa en tres pilares, todos ellos de naturaleza distinta al argumento de equidad que sustenta a SHS:
\begin{la}
    \item \textbf{Movilidad internacional del capital}. En una economía con libre circulación de capitales dentro de la Unión, gravar el ahorro y la inversión a los mismos tipos marginales elevados que el trabajo genera un fuerte incentivo a la deslocalización del capital hacia jurisdicciones más benignas, algo que no ocurre con el trabajo, mucho menos móvil. El argumento es que la elasticidad de la base imponible del capital a los tipos impositivos es mucho mayor que la del trabajo, lo que hace más eficiente gravar el capital a tipos más bajos, en términos de minimizar la pérdida de bienestar por unidad de recaudación;\footnote{Esta es la intuición detrás del resultado teórico de CHAMLEY-JUDD sobre la conveniencia de un tipo cero sobre el capital en el largo plazo en modelos de crecimiento óptimo, aunque ese resultado es aun objeto de fuerte controversia técnica.
    }
    \item \textbf{Integración con el Impuesto sobre Sociedades}. Un tipo proporcional sobre el capital facilita coordinar la tributación societaria (tipo fijo en el IS) con la tributación personal del accionista, evitando algunos problemas de doble imposición económica; y,
    \item \textbf{Simplicidad administrativa}. Un tipo único o casi único sobre las rentas del capital reduce el margen de arbitraje entre distintas formas de ahorro, frente a una escala progresiva que trataría de forma muy distinta a un depósito bancario y a una cartera de acciones según el tramo marginal del contribuyente.
\end{la}

Frente a esto, la crítica de equidad señala dos problemas: (i) \textbf{inequidad horizontal}, porque dos contribuyentes con idéntica renta total tributan de forma distinta según la composición de esa renta; y, (ii) \textbf{pérdida de progresividad efectiva en la cúspide} de la distribución, porque la renta del capital está mucho más concentrada en los deciles superiores de renta que la del trabajo, de modo que gravarla a tipos más bajos reduce la progresividad global del sistema precisamente donde más recaudación potencial hay.\footnote{El propio Informe MIRRLEES dedica un capítulo entero a esta tensión, concluyendo que no existe una solución de diseño que resuelva simultáneamente la neutralidad entre fuentes y la protección frente a la movilidad internacional del capital, y que la elección es, en última instancia, una decisión de política, no una cuestión técnica resuelta.
}

\subsubsection{Arbitraje de calificación: ¿cómo escoger mi fuente de renta?}
La existencia de una brecha significativa entre el tipo marginal máximo de la base general y el de la base del ahorro genera un incentivo racional a recalificar renta del trabajo como renta del capital para beneficiarse del tipo más bajo. 

El caso paradigmático, es el de los socios profesionales: un profesional (médico, abogado, arquitecto) que presta sus servicios a través de una sociedad de la que es socio mayoritario puede, en principio, optar por remunerarse vía salario (tributando en la base general, a tipo progresivo) o vía dividendo (tributando en la base del ahorro, a tipo inferior), a pesar de que económicamente la renta obtenida retribuye, en ambos casos, su trabajo personal. 

Para cerrar este arbitraje, la Ley 26/2014 introdujo un tercer párrafo en el art. $27.1$ LIRPF, según el cual los rendimientos obtenidos por el contribuyente procedentes de una entidad en cuyo capital participe, derivados de actividades profesionales, se califican como rendimientos de actividad económica cuando el contribuyente esté incluido en el Régimen Especial de Trabajadores Autónomos (RETA) o en una mutualidad de previsión social alternativa.\footnote{Esta regla ha generado abundante litigiosidad reciente. El Tribunal Supremo ha aclarado, que la calificación como actividad económica no exige el alta formal en el RETA, sino que basta con que el contribuyente esté obligado legalmente a cotizar en ese régimen por control societario efectivo, primando así la realidad económica sobre el formalismo administrativo. Se trata, en definitiva, de un ejemplo claro de cómo una divergencia estructural genera, casi mecánicamente, comportamiento elusivo, y de cómo el legislador y los tribunales deben perseguirlo con normas y doctrina específicas.
}

\subsubsection{Rentas no gravadas: ¿se grava realmente la capacidad de pago?}
Por otro lado, la imputación de rentas y la percepción de pagos en especie son también un caso claro de divergencias por fuente en el tratamiento que se le da en España. Se trata de renta que SHS consideraría plenamente gravable, ya que constituye poder de consumo real, pero que el IRPF grava de forma parcial según cuál sea su naturaleza concreta.

\paragraph{Imputación de rentas inmobiliarias: gravamen asimétrico}
Por la titularidad de inmuebles urbanos no arrendados ni afectos a actividades económicas, el contribuyente debe imputarse una renta ficticia equivalente al $2\%$ del valor catastral, en proporción al número de días de disponibilidad del inmueble en el ejercicio ($1,1\%$ si el valor catastral ha sido revisado en los últimos diez años).

Sin embargo, la propia norma excluye expresamente de esta imputación a la vivienda habitual. La consecuencia es una asimetría flagrante desde el punto de vista SHS: el propietario de una segunda residencia tributa por el ahorro de alquiler que le proporciona disponer de ella, mientras que el propietario que habita su única vivienda no tributa por ella en absoluto, aunque obtiene exactamente el mismo tipo de renta en especie: el valor de uso de un activo propio. 

Esta es, en la literatura de finanzas públicas, la divergencia más citada tras la de las plusvalías no realizadas: el consumo de vivienda en propiedad es, con diferencia, la mayor partida de renta imputada que un sistema SHS puro debería gravar y que casi ningún sistema tributario real grava de forma completa por razones más políticas y sociales que técnicas.

\paragraph{En esepcie: ¿cómo escoge la forma de mi remuneración?}
El caso paradigmático es la definición del art. $42.1$ LIRPF. Este define como renta en especie la utilización, consumo u obtención, para fines particulares, de bienes, derechos o servicios de forma gratuita o por precio inferior al de mercado, y como regla general las valora e integra en la base general como mayor rendimiento del trabajo.

Ahora bien, el mismo artículo exime de tributación una lista tasada de retribuciones en especie: cheques-restaurante o fórmulas de comedor de empresa (hasta el límite reglamentario diario), primas de seguro de enfermedad hasta $500€$ anuales por beneficiario ($1.500$ euros en caso de discapacidad), gastos de formación exigidos por el puesto de trabajo, guarderías, o la compensación por kilometraje en vehículo propio hasta $0,19€$/km.

Formalmente son exenciones, pero ilustran, desde el lado del trabajo, el mismo fenómeno que la vivienda habitual ilustra desde el lado del capital: hay renta real, con capacidad económica indiscutible, que el legislador decide no gravar selectivamente según cuál sea su forma de percepción. De nuevo, confirma que la ruptura de la comprehensividad de fuente en el IRPF español no se limita a la frontera trabajo/capital, sino que atraviesa también la frontera dinerario/especie.

\subsection{Debates abiertos: ¿Es mejorable la arquitectura del impuesto?}
En definitiva, vemos que existe una clara divergencia entre la base imponible del impuesto y el benchmark óptimo desde el punto de vista teórico. Estas no se reducen a las expuestas, existen otros debates abiertos en torno a otras figuras en los que hay renta, en sentido SHS, que el legislador podría integrar en el IRPF con toda naturalidad técnica, y decide no hacerlo por razones ajenas a la lógica del propio impuesto.

\subsubsection{Transmisiones gratuitas: Mortis causa}
El art. $6.4$ LIRPF excluye de sujeción al IRPF las rentas sujetas al Impuesto sobre Sucesiones y Donaciones (ISD), y el art. 33.3.b) LIRPF declara no sujetas las ganancias patrimoniales puestas de manifiesto por transmisiones mortis causa. 

Desde la perspectiva SHS, esta exclusión es difícil de justificar en términos puramente analíticos: recibir una herencia o una donación es, sin ningún género de dudas, un incremento del poder de disposición económica del receptor originado por la liberalidad ajena. ¿Por qué, entonces, prácticamente ningún sistema comparado integra las adquisiciones gratuitas en el impuesto sobre la renta, prefiriendo un tributo separado?\footnote{Combinación con la transmisión mortis causa: de diferimiento a exención definitiva. El art. 33.3.b) LIRPF establece que no existe ganancia o pérdida patrimonial con ocasión de transmisiones lucrativas por causa de muerte del contribuyente —la llamada coloquialmente "plusvalía del muerto"—. La consecuencia es que toda la revalorización acumulada por el causante durante su vida, y nunca realizada, desaparece definitivamente del ámbito del IRPF: no tributa en su cabeza ni, en principio, en la del heredero, que además recibe el bien con un nuevo valor de adquisición (de mercado) que le servirá de referencia para una futura venta. Esto convierte lo que en vida era un mero diferimiento del impuesto en una exención definitiva: la renta SHS generada nunca es objeto de gravamen alguno en el IRPF, ni por el causante ni por el heredero. Es exactamente el mismo fenómeno conocido en la literatura anglosajona como stepped-up basis at death o, coloquialmente, angel of death loophole, y ha sido objeto de propuestas legislativas recientes en Estados Unidos (WYDEN, 2021) precisamente a raíz de investigaciones periodísticas sobre los tipos efectivos extraordinariamente bajos que soportan las grandes fortunas cuya riqueza crece principalmente vía revalorización de activos no realizados.

No obstante, el legislador español sí ha detectado y corregido un uso abusivo concreto de esta combinación. La Ley 11/2021, de medidas de prevención y lucha contra el fraude, modificó el art. $36$ LIRPF para excluir del beneficio de la no sujeción a quienes, mediante pactos sucesorios con efectos de presente (presente en algunos derechos forales: gallego, catalán), transmitan el bien en vida del causante pero con calificación mortis causa, y posteriormente lo vendan dentro de los cinco años siguientes: en ese caso, el adquirente se subroga en el valor y fecha de adquisición del causante, en vez de beneficiarse del valor actualizado. 
} Las razones que ofrece la doctrina son, en su mayoría, ajenas a la lógica interna de SHS:
\begin{la}
    \item \textbf{Antigëdad}. En primer lugar, el gravamen de las transmisiones mortis causa e inter vivos es, en España y en la mayoría de ordenamientos continentales, mucho más antiguo que el impuesto sintético sobre la renta. Tiene su origen en los derechos reales de la Hacienda medieval sobre las transmisiones patrimoniales, y se ha desarrollado como una rama autónoma del sistema tributario, con reglas propias de valoración de patrimonios y de parentesco, difícilmente trasladables sin más al esquema anual y de flujo del IRPF;
    \item \textbf{Focos distintos: patrimonio vs capacidad económica}. En segundo lugar, un impuesto sobre transmisiones patrimoniales pone el \textbf{foco en el patrimonio} que se transmite; con tarifas que varían según el grado de parentesco, atendiendo a una lógica de solidaridad familiar. Por el contrario, un IRPF integrado pondría el \textbf{foco en la capacidad económica} del receptor considerada aisladamente. Por tanto, reprdocuen dos lógicas de equidad distintas que no siempre coinciden en sus resultados; y,
    \item \textbf{Costes políticos}. Por último, y probablemente más importante, el ISD es, junto con el Impuesto sobre el Patrimonio, uno de los tributos con menor aceptación social en España (impuesto de la muerte) y su fuerte cesión normativa a las Comunidades Autónomas ha generado una dispersión de tipos efectivos extraordinaria entre territorios, con bonificaciones que en algunas Comunidades reducen la tributación efectiva entre parientes directos a niveles casi simbólicos, mientras que en otras se mantiene una carga sustancial. Esta dispersión es, en sí misma, un argumento adicional a favor de mantener el ISD separado: integrar las herencias en un IRPF de base estatal con cesión normativa parcial del $50\%$ generaría fricciones de reparto competencial mucho más complejas que las que ya genera el propio IRPF.
\end{la}

No obstante, conviene mencionar que el ISD no es un sustituto neutral de una hipotética integración en el IRPF.\footnote{Por ejemplo, la reducción del 95\% en la base imponible del ISD por transmisión de empresa familiar, reducción que algunas Comunidades Autónomas han elevado hasta el 99\%, es un ejemplo interesante porque, aunque formalmente pertenece al ISD y no al IRPF, ilustra el mismo patrón de fondo que veremos más adelante: una renta claramente gravable en términos SHS queda prácticamente exenta por una decisión de política económica.
}

Otro ejemplo, más ligero en apariencia pero analíticamente idéntico serían los premios de loterías y apuestas. Una ganancia de esta naturaleza es, en términos SHS, indistinguible de cualquier otra ganancia patrimonial y, sin embargo, el legislador ha optado por sacarla del régimen general de ganancias patrimoniales del IRPF y darle un tratamiento autónomo, más favorable en la franja baja gracias al mínimo exento de $40.000$ euros. Como podemos suponer, la explicación no es técnica, sino de aceptabilidad social. La Lotería de Navidad es una institución cultural profundamente arraigada en España, y un régimen de tributación percibido como excesivamente gravoso sobre habría tenido un coste político desproporcionado respecto de la recaudación en juego.

\subsubsection{Neutralidad intertemporal estricta: ¿es neutral el IRPF?}
El segundo gran debate abierto no cuestiona el qué se grava (la fuente), sino el cuándo. SHS grava la renta en el momento en que se genera, con independencia de que el contribuyente decida consumirla de inmediato o ahorrarla para consumo futuro. 

Esta aparente neutralidad esconde, sin embargo, un problema que la literatura de finanzas públicas ha señalado desde hace décadas: un impuesto sobre la renta que grava tanto el principal ahorrado como el rendimiento que ese principal genera en el futuro grava dos veces la decisión de posponer el consumo, mientras que quien consume de inmediato solo tributa una vez. 

Es el argumento central de KALDOR (1955): \textbf{si dos contribuyentes obtienen la misma renta pero uno la consume íntegramente y el otro la ahorra e invierte, bajo un impuesto de renta comprehensivo el segundo acaba pagando más impuesto en valor presente que el primero, por el simple hecho de haber ahorrado}; lo que constituye, según esta tesis, una falta de neutralidad intertemporal que penaliza el ahorro frente al consumo. La alternativa propuesta por KALDOR  sería un impuesto sobre el gasto personal, que grava la renta menos lo ahorrado, es decir, solo lo efectivamente consumido. Esta alternativa fue retomada y desarrollada técnicamente en el Informe MEADE, que sigue siendo la referencia clásica en el debate renta-consumo como base alternativa de imposición directa.

\paragraph{Materialización en España: Previsión Social Complementaria}
España no ha adoptado un impuesto sobre el gasto, pero sí ha introducido, para una categoría concreta de ahorro, un tratamiento que se aproxima al ideal de neutralidad intertemporal defendido por esa corriente: la previsión social complementaria. 

El régimen fiscal de las aportaciones a planes de pensiones y sistemas equivalentes sigue el modelo conocido en la literatura comparada como EET (Exento-Exento-Tributado): las aportaciones reducen la base imponible en el momento de realizarse, el rendimiento generado dentro del fondo no tributa año a año mientras permanece invertido, y solo se somete a gravamen como rendimiento del trabajo en el momento del rescate. Como vemos, este esquema replica, con bastante fidelidad, el resultado de un impuesto sobre el gasto: lo que se grava, en términos de valor actual, es únicamente la renta que finalmente se consume, no la que se destina a ahorro previsional. 

Ahora bien, este tratamiento no se generaliza al resto del ahorro. El ahorro ordinario (depósitos, fondos de inversión, carteras de valores) sigue un esquema mucho más próximo al TTE que al EET: se invierte con renta ya gravada (no hay reducción en el momento de aportar), tributa anualmente por los rendimientos del capital mobiliario efectivamente percibidos (dividendos, intereses), y tributa de nuevo, al criterio de realización ya analizado, por la ganancia patrimonial en el momento de la venta. El resultado es una falta de neutralidad entre instrumentos de ahorro, no solo entre trabajo y capital: el IRPF trata de forma sustancialmente más favorable, en términos de valor actual, al ahorro canalizado a través de un plan de pensiones que al ahorro canalizado a través de una cuenta de valores ordinaria, sin que exista una justificación de capacidad económica para esa diferencia —solo una decisión de política de fomento del ahorro finalista para la jubilación, que es, de nuevo, un argumento de política económica y no de coherencia con el benchmark SHS.

\subsection{Balance: ¿qué tipo de impuesto es, en rigor, el IRPF?}
A lo largo del apartado hemos identificado tres grandes bloques de divergencia estructurales respecto al óptimo teórico: (i) una divergencia temporal; (ii) una divergencia por fuente; y, (iii) dos debates abiertos sobre la propia arquitectura del impuesto, referidos a qué figuras quedan fuera de su ámbito material (transmisiones gratuitas, premios de juego) y a la falta de neutralidad intertemporal entre distintos instrumentos de ahorro.

En términos prácticos resulta imposible llevar a cabo un diseño conforme al criterio de Renta SHS. Por eso, la literatura de finanzas públicas comparada distingue, de forma estilizada, varios modelos de imposición personal sobre la renta según su grado de fidelidad al ideal SHS: (i) el \textbf{impuesto cedular} o de producto, que grava cada fuente de renta de forma aislada, sin acumulación ni progresividad conjunta; (ii) el \textbf{impuesto sintético puro}, fidelidad máxima a la comprehensividad de fuente, con una única escala progresiva sobre la totalidad de la renta, con independencia de su origen; (ii) el \textbf{Dual Income Tax} (DIT) puro, de origen nórdico, que combina una escala progresiva sobre las rentas del trabajo con un tipo proporcional único, y generalmente bajo, sobre la totalidad de las rentas del capital; y, (iv) una categoría intermedia que la literatura reciente denomina \textbf{semi-dual income tax}, para describir sistemas que introducen una base separada para (parte de) las rentas del capital, pero sin llevar la lógica dual hasta sus últimas consecuencias. 

Como savemos, España se encuedra, a priori Entonces, lo que ahora conviene señalar es el tipo de benchmark que establecemos al analizar los gastos fiscales en España y su encuadre dentro de la literatura existente. ¿Qué tipo de impuesto es el IRPF?

Como sabemos, el IRPF español, al igual que el resto de sistemas impositivos reales, rompe deliberadamente la exigencia de gravamen SHS: clasifica la renta en dos bases con tarifas distintas.\footnote{La base general tributa según una escala progresiva por tramos que combina tramo estatal y tramo autonómico; mientras que la base del ahorro tributa según una escala separada, también progresiva por tramos desde la reforma de 2010, pero con tipos sensiblemente inferiores.
} 

\subsubsection{Modelo semi-dual}
La doctrina especializada en fiscalidad comparada, y la propia literatura académica española sobre el IRPF dual, desde los primeros trabajos empíricos de CABRÉ (2001) hasta los estudios de DÍAZ CARO, ONRUBIA y PÉREZ MAYO (2013), coincide en que el IRPF español, pese a estructurarse formalmente en dos bases no se ajusta a un modelo dual puro, y lo caracteriza más propiamente como un modelo semi-dual. 

Son tres las razones concretas que sustentan esta calificación, y conviene dominarlas con precisión. En primer lugar, la \textbf{base del ahorro no agota todas las rentas del capital}. A diferencia del modelo nórdico puro, donde la totalidad de la renta del capital tributa al tipo proporcional reducido, en España los rendimientos del capital inmobiliario (alquileres) y determinadas ganancias patrimoniales que no derivan de una transmisión (por ejemplo, ciertas subvenciones vinculadas a la vivienda habitual) permanecen integrados en la base general, sujetos a la escala progresiva. Esto significa que la frontera dual española no coincide con la frontera trabajo/capital, sino más bien con la frontera capital financiero/capital no financiero —una distinción de diseño mucho más fina, y también más arbitraria desde el punto de vista de la neutralidad, que la del modelo nórdico. En segundo lugar, la \textbf{escala de la base del ahorro no es proporcional}, sino progresiva por tramos. Como vimos, la base del ahorro no tributa a un tipo único, sino según una escala de varios tramos, lo que atempera la brecha con la base general, pero aleja aún más al modelo español del ideal dual puro de tipo único sobre el capital, que fue precisamente el diseño originario. Por último, la \textbf{ausencia de integración plena con el Impuesto sobre Sociedades}. El modelo dual nórdico busca, como uno de sus objetivos de diseño, que el tipo proporcional sobre el capital personal coincida o se aproxime al tipo del impuesto sobre sociedades, minimizando el arbitraje entre tributar beneficios en sede societaria o personal. En España, el tipo general del Impuesto sobre Sociedades ($25\%$) y los tramos de la base del ahorro ($19-30\%$) no están alineados de forma sistemática, lo que sigue generando incentivos a la planificación fiscal entre ambas figuras.

La consecuencia práctica de esta caracterización es que España combina lo peor y lo mejor de ambos mundos, según el criterio con que se mire. Por un lado, conserva buena parte de la complejidad de fronteras del modelo sintético (hay que calificar cada renta por su origen exacto para saber a qué base va). Por otro lado, no presenta toda la simplicidad ni toda la eficiencia recaudatoria sobre el capital móvil que promete el modelo dual puro.


\clearpage
\section{Divergencias intencionales: Gastos Fiscales}
\puntosclave{
    
}   

En el apartado anterior logramos definir un concepto clave: el sistema tributario de referencia real del IRPF español, la base sobre la que opera la Hacienda Pública en su Memoria de Beneficios Fiscales. A partir de ahora, valoremos a partir de este benchmark semi-dual ya caracterizado las desviaciones intencionales introducidas por el legislador con fines de política económica y social: los gastos fiscales propiamente dichos.

Ahora bien, no toda minoración del impuesto es lo mismo solo porque en la conversación cotidiana llamemos "deducción" a cualquier ventaja fiscal. Jurídica y económicamente son instrumentos distintos, actúan en momentos distintos de la liquidación, y —esto es lo importante— producen efectos distributivos distintos aunque el "importe nominal" del beneficio parezca idéntico. Si memorizas el catálogo del bloque 2 sin haber fijado antes esta tipología, vas a describir los beneficios fiscales sin poder valorarlos, que es exactamente lo que un tribunal de TCEE espera que sepas hacer.

\subsection{Forma jurídica: ¿Por qué es importante?}
Antes de entrar en el catálogo, hay que dejar claro por qué la forma jurídica de un beneficio fiscal merece un tratamiento propio y no una mención de pasada. Hay tres razones, y las tres tienen consecuencia práctica directa para cualquier análisis posterior del catálogo previsto actualmente.

\subsubsection{Instrumentos: ¿tienen los mismos efectos distributivos?}
En primer lugar, es una \textbf{exigencia legal}, no una elaboración doctrinal externa. El art. 8.d) LGT somete a reserva de ley el establecimiento, modificación, supresión y prórroga de las exenciones, reducciones, bonificaciones, deducciones y demás beneficios o incentivos fiscales. La LIRPF, a su vez, las ubica en tramos distintos y sucesivos de su propio articulado, lo que confirma que la distinción no es un artificio sino la propia arquitectura interna de la norma que estamos estudiando. Segunda, porque de esa forma jurídica depende, analíticamente, \textbf{quién se beneficia y cuánto}, cuestión sobre la que hablaremos a continuación. Tercera, porque el propio \textbf{texto constitucional} da pie a la pregunta. El art. $31.1$ CE exige que el sistema tributario se inspire en los principios de igualdad y progresividad. Un beneficio fiscal mal diseñado puede erosionar silenciosamente esa progresividad constitucionalmente exigida, incluso cuando el legislador declara perseguir un objetivo redistributivo con la medida. Este es, en el fondo, el problema de fondo que subyace a todo el subapartado: la intención declarada de una norma y su efecto distributivo real pueden divergir por una simple decisión de técnica legislativa, y detectar esa divergencia es, precisamente, lo que se te va a pedir que sepas hacer.

\paragraph{Exención: Efecto marginal creciente} 
La renta exenta no llega siquiera a integrarse en los rendimientos íntegros: desaparece del cálculo antes de empezar. 

Su efecto distributivo es el más radical de todos los instrumentos, por dos razones acumulativas: primero, como excluye renta de una base que de otro modo tributaría de forma progresiva, su valor en euros ahorrados crece con el tramo marginal del beneficiario exactamente igual que una reducción (lo desarrollo con cifras a continuación); segundo, y esto es lo que la distingue de la reducción, al no aparecer siquiera como partida declarada, escapa incluso al control de rentas exentas con progresividad que si aplica a otras figuras, lo que dificulta más su seguimiento estadístico y su cuantificación en la Memoria de Beneficios Fiscales. 

El ejemplo de la exención de indemnizaciones por despido (hasta $180.000$ euros) ilustra bien el problema: como el importe exento es un límite absoluto en euros, no relativo a la retribución del trabajador, favorece proporcionalmente más a quien tenía un salario más elevado.

\paragraph{Reducción en la base imponible: Efecto distributivo creciente} 
Las reducciones en la base imponible minoran la base antes de aplicar la tarifa progresiva.

Aquí el efecto es cuantificable con precisión: si dos contribuyentes, A en el tramo marginal del $19\%$ y B en el tramo marginal del $45\%$, disfrutan cada uno de una reducción de $1.000$ euros, el ahorro fiscal real es de $190$ euros para A y de $450$ euros para B. Es decir, más del doble por la simple razón de que B tributaba, sobre ese tramo de renta, a un tipo más alto. 

Se trata del instrumento con el peor perfil distributivo de todo el catálogo, porque subsidia más, en términos absolutos, a quien menos lo necesita en cualquier lectura de capacidad económica. Este es, precisamente, el fundamento técnico de una de las críticas más repetidas en la literatura la reducción por aportaciones a planes de pensiones: al operar en base y no en cuota, beneficia desproporcionadamente a los tramos altos de renta, que son además los que menos necesitan un incentivo adicional para ahorrar, mientras que un asalariado de renta baja, con escasa o nula capacidad de ahorro, apenas puede aprovecharla.

\paragraph{Deducción en la cuota: Proporcionalmente neutro} 
Las deducciones en la cuota se restan directamente del impuesto ya calculado. 

Con la misma cifra del ejemplo anterior, una deducción de 190 euros afecta por igual tanto A como B, con independencia de su tramo marginal: es un instrumento proporcionalmente neutro, y si el importe de la deducción es fijo en vez de porcentual sobre una base variable, incluso levemente progresivo en términos relativos, porque representa una fracción mayor de la cuota total de quien tiene una cuota más baja.

\paragraph{Deducciones rembolsables o anticipables: Genuinamente progresiva}
Las deducciones reembolsables o anticipales traen la lógica anterior a su extremo: minoran la cuota diferencial, se cobran mensualmente por anticipado, y si la cuota resulta insuficiente para absorberlas, o el contribuyente ni siquiera está obligado a declarar, la Hacienda Pública abona igualmente el importe. Por ejemplo, por maternidad y familia numerosa/discapacidad.

Es el único instrumento del catálogo cuyo efecto distributivo no depende en absoluto de la renta del beneficiario: un contribuyente sin ninguna cuota que pagar recibe exactamente el mismo importe que uno con cuota sobrada, lo que la convierte, de las cinco figuras analizadas, en la única genuinamente progresiva en sentido estricto —beneficia relativamente más a quien menos renta tiene—. 

Pero esta misma característica genera un problema de implementación real y documentado: al convertir una deducción fiscal en un pago periódico anticipado, la Administración asume el riesgo de cobros indebidos cuando las circunstancias del beneficiario cambian a lo largo del año (pérdida de la condición de familia numerosa, superación de umbrales de discapacidad, etc.), lo que obliga a un procedimiento específico de regularización mediante el Modelo 122, con el consiguiente riesgo de que el contribuyente deba devolver cantidades ya gastadas. Una fricción administrativa que no existe en ninguna de las figuras anteriores, precisamente porque estas no implican flujo de caja anticipado.

\paragraph{Diferimento: Regresivo (distribución personal de la riqueza)}
La figura del diferimento ya fue analizada a propósito del criterio de realización: no elimina el gravamen, lo traslada en el tiempo. 

Para no redesarrollarlo aquí, conviene fijar su efecto distributivo propio, que no tratamos entonces: un diferimiento beneficia más a quien puede permitirse no vender el activo revalorizado mientras espera un momento fiscalmente más favorable (lo que exige, por definición, disponer de liquidez alternativa suficiente para no depender de esa venta), de modo que su ventaja se concentra, de forma sistemática, en los patrimonios más elevados y más diversificados.

\paragraph{Tipo de gravamen reducido: Depende}
El tipo de gravamen reducido tiene un efecto distributivo que es el más matizado de los seis: (a) si el tipo reducido se aplica desde el primer euro, es regresivo en el sentido de que reduce más, en términos relativos, la carga de quien más obtiene; pero, (b) si se combina con un mínimo exento antes de aplicar ese tipo reducido (como ocurre con el gravamen sobre loterías), el resultado neto puede ser incluso progresivo en la franja baja y solo proporcional a partir de cierto umbral, lo que obliga a analizar cada caso en concreto y no dar por supuesto un signo distributivo único para esta categoría.

\subsubsection{Institucional: ¿cómo actúa el Estado y cómo esto influye en el debate?}
Además, existe un problema adicional que cabe valorar adecuadamente: el elemento institucional, que tiene un origen histórico preciso que conviene conocer. Fue el propio SURREY quien, tras años de insistir desde el Tesoro estadounidense en que los beneficios fiscales eran, en sustancia, gasto público disfrazado, logró que esa idea se trasladara a la práctica presupuestaria: la Congressional Budget and Impoundment Control Act de 1974, en Estados Unidos, obligó por primera vez a que el presupuesto federal incluyera un presupuesto de gastos fiscales junto al presupuesto de gasto directo. 

La lógica de fondo era exactamente la que se planteó antes: si un beneficio fiscal no pasa por el debate parlamentario anual de la partida de gasto correspondiente, escapa al escrutinio político ordinario, y esa invisibilidad tiene consecuencias reales sobre qué políticas se mantienen y cuáles se cuestionan.

Esta intuición fue desarrollada después, en el ámbito académico, por C. HOWARD (\textit{The Hidden Welfare State}, 1997)), un trabajo que sostiene una tesis provocadora y muy citada en la literatura de política social comparada: buena parte de la redistribución efectiva que realiza el Estado se produce a través de beneficios fiscales y no a través de programas de gasto social directo, con la particularidad de que estos beneficios fiscales generan menos controversia política que un cheque o una prestación equivalente, a pesar de perseguir exactamente el mismo objetivo redistributivo, porque se perciben culturalmente como una rebaja de impuestos y no como una ayuda pública.

HOWARD documenta, además, que este Estado del bienestar oculto tiende a beneficiar más a las clases medias y medias-altas que a los colectivos de renta más baja, precisamente por el problema de regresividad de las reducciones en base que acabamos de analizar.\footnote{El caso paradigmático en la literatura estadounidense es la deducción por intereses hipotecarios, que durante décadas benefició desproporcionadamente a los propietarios de las viviendas más caras, en comparación con un hipotético programa de ayuda directa a la vivienda de igual coste presupuestario pero mejor focalizado.
} España no ha sido ajena a este problema, y la creación del Presupuesto de Beneficios Fiscales responde exactamente a la misma lógica que la reforma estadounidense de 1974: hacer visible, mediante un documento presupuestario específico, un gasto que de otro modo quedaría oculto en el articulado de cada tributo. Pero conviene no sobrevalorar el alcance real de esta transparencia formal: la propia doctrina hacendística ha señalado que la Memoria de Beneficios Fiscales, aunque técnicamente completa, funciona en la práctica como un anexo informativo al proyecto de Presupuestos y no recibe, en el debate parlamentario de las Cortes Generales, un escrutinio comparable, partida por partida, al que reciben los programas de gasto directo, que deben defenderse y justificarse cada año ante la Comisión de Presupuestos. 

Esta es la razón por la que una corriente de reforma, con eco en recomendaciones de la OCDE sobre buenas prácticas presupuestarias, propone introducir cláusulas de caducidad automática (sunset clauses) para los beneficios fiscales de mayor cuantía; de modo que, transcurrido un plazo determinado, decaigan salvo renovación legislativa expresa, obligando así a una reevaluación periódica que hoy no existe. Frente a esta propuesta, el régimen español actual, en el que un beneficio fiscal, una vez incorporado a la Ley, permanece indefinidamente vigente hasta que una reforma expresa lo derogue, con independencia de si sigue cumpliendo o no el objetivo que originalmente lo justificó.

\subsection{Catálogo de gastos fiscales: ¿Qué beneficios fiscales contempla el IRPF?}
Ahora que conocemos las figuras y efectos de los diferentes instrumentos fiscales, pasemos a analizar algunos de los que efectivamente se aplican en España. 

Frente a la tentación de recorrer el articulado de la LIRPF de forma correlativa, agrupemos los beneficios fiscales según el objetivo de política pública que el legislador declara perseguir: vivienda, familia, ahorro, emprendimiento, mecenazgo y empleo. Esta agrupación no es solo una comodidad expositiva: cada finalidad tiene su propia literatura de evaluación de políticas públicas, sus propios defensores y críticos, y sus propios problemas de diseño recurrentes, de modo que entender un bloque te permite razonar sobre cualquier medida nueva que se introduzca dentro de esa misma finalidad, aunque no la hayas estudiado en concreto —que es exactamente la habilidad que un ejercicio oral pretende comprobar.

\subsubsection{Vivienda: Deducción}
El acceso a la vivienda en propiedad ha sido, desde los años noventa, uno de los objetivos de política social más constantes de la fiscalidad española, apoyado en un argumento de estabilidad patrimonial y de ahorro forzoso para las familias, además de en la función macroeconómica del sector de la construcción como motor de actividad y empleo.

La deducción por adquisición de vivienda habitual, fue derogada en 2013, manteniendo un régimen transitorio para quienes hubieran adquirido su vivienda habitual, o satisfecho cantidades para su construcción, antes de esa fecha, incluyendo expresamente la deducibilidad de las cantidades destinadas a la cancelación del préstamo hipotecario correspondiente. La deducción por alquiler de vivienda habitual del inquilino siguió el mismo camino, derogada también desde 2015 salvo régimen transitorio para contratos anteriores. 

En su lugar, el legislador ha desplazado el incentivo hacia el propietario-arrendador: la Ley 12/2023, por el derecho a la vivienda, escalona las reducciones en el rendimiento neto del capital inmobiliario por alquiler de vivienda habitual: hasta el $90\%$ en zonas de mercado tensionado con rebaja de renta, $70\%$ para inquilinos jóvenes o rehabilitaciones recientes, $60\%$ general en rehabilitación, y un $50\%$ de suelo si no concurre ninguna de las anteriores, trasladando así el incentivo fiscal de comprar a poner en alquiler en condiciones sociales.\footnote{Más recientemente, el Real Decreto-ley 19/2021, en el marco del Plan de Recuperación, Transformación y Resiliencia financiado con fondos Next Generation EU, introdujo tres deducciones escalonadas por obras de mejora de la eficiencia energética: del $20\%$ (reducción de demanda de calefacción/refrigeración, hasta 5.000 euros de base), del $40\%$ (mejora del consumo de energía primaria no renovable, hasta 7.500 euros) y del $60\%$ (rehabilitación energética de edificios completos, hasta 15.000 euros), condicionadas a la acreditación mediante certificado energético antes y después de la obra. Por último, dos exenciones de ganancias patrimoniales cierran el bloque: para mayores de 65 años, la venta de la vivienda habitual queda exenta sin necesidad de reinversión, y la venta de cualquier otro activo patrimonial queda exenta si el importe se reinvierte en una renta vitalicia en un plazo de seis meses, con el límite de 240.000 euros.
}

Este es, de todos los bloques del catálogo, el que cuenta con una literatura crítica más consolidada, y el argumento central es el de la capitalización del beneficio fiscal en el precio. POTERBA (1984), en un trabajo ya clásico sobre la fiscalidad de la vivienda en propiedad, mostró con un modelo de mercado de activos que, cuando la oferta de vivienda es poco elástica en el corto y medio plazo, un incentivo fiscal a la demanda (una deducción por compra) tiende a trasladarse, total o parcialmente, al precio de venta del inmueble, de modo que buena parte del beneficio termina en manos del vendedor y no del comprador al que se dirigía la política. 

Esta crítica, sumada al argumento distributivo propia de las deducciones, fue uno de los fundamentos técnicos que el Informe LAGARES (2014) invocó para justificar y consolidar su supresión.\footnote{El régimen transitorio, sin embargo, sigue vigente más de una década después de la derogación formal, y su coste recaudatorio, aunque decreciente año a año a medida que se amortizan las hipotecas afectadas, todavía representaba una de las principales partidas de la Memoria de Beneficios Fiscales del IRPF en los últimos ejercicios —un ejemplo perfecto, para cerrar este bloque, del problema de persistencia normativa que analizamos: una medida derogada en el papel puede seguir teniendo un coste presupuestario relevante durante décadas.
}


\subsubsection{Familia, natalidad y discapacidad: Deducciones anticipadas}
Frente al argumento puramente estructural del mínimo por descendientes o discapacidad que fijamos como divergencia estructural, centrémonos en las medidas cuya finalidad declarada es explícitamente de fomento demográfico y de apoyo a la conciliación, en un país con una de las tasas de natalidad más bajas de la Unión Europea. Ya analizamos la mecánica distributiva de las deducciones, bien sean por maternidad y/o por familia numerosa o discapacidad de ascendientes/descendientes. Así que no repetiremos su mecánica de abono anticipado, ya vista, sino su encaje dentro de esta finalidad concreta.

El debate de fondo en este bloque no es tanto de eficacia distributiva sino de eficacia demográfica: ¿incentivan realmente estas deducciones la natalidad, o simplemente transfieren renta a familias que habrían tenido hijos de todos modos? La literatura internacional sobre políticas de natalidad (Francia o los países nórdicos, que llevan décadas experimentando con incentivos fiscales y prestaciones familiares) tiende a mostrar elasticidades modestas de la fecundidad respecto de estos incentivos monetarios, sugiriendo que factores como la disponibilidad de plazas de guardería pública o la flexibilidad laboral pesan más que el importe de la deducción en la decisión de tener hijos, lo cual no invalida la medida como instrumento redistributivo hacia las familias con hijos, pero sí cuestiona su eficacia como palanca demográfica en sentido estricto.

Conviene mencionar, como dato de conjunto del sistema español de protección a la familia, que estas deducciones del IRPF coexisten desde 2020 con el Ingreso Mínimo Vital, prestación no contributiva de la Seguridad Social gestionada al margen del sistema tributario. Es decir, el propio legislador español, optó por un instrumento de gasto directo y no por ampliar el catálogo de deducciones del IRPF, en una decisión que puede leerse como un reconocimiento implícito de los límites de la vía fiscal para llegar a los hogares sin ingresos suficientes para generar cuota alguna que minorar por vía de deducción no anticipable.

\subsubsection{Ahorro y previsión social: Deducción sobre la base}
Como vimos en el bloque anterior, uno de los objetivos perseguidos por esta figura impositiva es fomentar el ahorro finalista para la jubilación complementaria, en un sistema de pensiones públicas de reparto sometido a presión demográfica.

De nuevo, omitiremos la mecánica distributiva de estas medidas EET ya desarrollada, y nos centraremos en señalar el dato que corresponde a este bloque y no al anterior: la Memoria de Beneficios Fiscales del Estado sí computa esta reducción como gasto fiscal cuantificado, con un coste que se ha reducido sensiblemente en los últimos ejercicios a medida que el legislador ha ido recortando el límite máximo de aportación deducible (de $8.000$ a $1.500$ anuales para aportaciones individuales desde la Ley de Presupuestos de 2021, compensado parcialmente con la ampliación del límite para aportaciones empresariales hasta 8.500 euros).

El problema de diseño de este bloque es el mismo que ya fijamos: al operar como reducción en base y no como deducción en cuota, beneficia desproporcionadamente a los tramos de renta alta, precisamente los que menos necesitan un incentivo adicional para ahorrar y los que con mayor probabilidad ahorrarían de todos modos sin el beneficio fiscal, lo que agrava el exceso de gravamen de la figura: gasto fiscal que no genera ahorro adicional, sino que simplemente subvenciona un ahorro que ya se iba a producir.\footnote{La evidencia empírica más citada sobre esta cuestión no proviene de España sino de Estados Unidos: DUFLO et al. (2006), en un experimento de campo realizado junto con la empresa de asesoría fiscal H\&R Block, mostraron que sustituir un incentivo articulado como deducción en base por uno articulado como crédito de igualación multiplicaba muy significativamente la tasa de apertura de cuentas de ahorro para la jubilación entre los hogares de renta baja y media, frente al diseño tradicional de reducción en base, que apenas modificaba su comportamiento. Es la confirmación empírica, con datos de campo y no solo con aritmética de tramos marginales, de la tesis distributiva que desarrollamos como marco teórico: la forma jurídica del incentivo determina no solo quién se beneficia, sino si el incentivo cumple o no su función de cambiar realmente el comportamiento del contribuyente.
}

\subsubsection{Emprendimiento e inversión empresarial: Deducción sobre cuota íntegra}
La idea general de este tipo de gastos fiscales es fomentar la financiación privada de proyectos empresariales en fase temprana, allí donde el crédito bancario tradicional resulta insuficiente o inaccesible por el riesgo asociado a empresas sin trayectoria.

Para el caso español, la deducción por inversión en empresas de nueva o reciente creación, permite deducir en la cuota íntegra estatal el $50\%$ de las cantidades invertidas en la suscripción de acciones o participaciones, con una base máxima de deducción de $100.000€$ anuales. La propia Ley de Startups introduce, además, un régimen reforzado para las empresas certificadas como emergentes por ENISA (Empresa Nacional de Innovación), con excepciones específicas para los socios fundadores (como la dispensa del límite de participación del 40\%).

El debate de fondo aquí es doble. Por un lado, el de eficacia: ¿genera este tipo de incentivo fiscal inversión adicional en el ecosistema emprendedor, o simplemente subvenciona inversiones que ya se iban a producir por otros motivos (rentabilidad esperada, diversificación de cartera)? KERR, LERNER y SCHOAR (2014), en un estudio de referencia sobre financiación en Estados Unidos, documentan que el acceso a financiación temprana tiene efectos causales positivos y significativos sobre la supervivencia y el crecimiento de las start-ups financiadas, lo que da cierto respaldo empírico a la lógica de fomento, aunque no aísla específicamente el efecto del incentivo fiscal frente al de la financiación en sí misma. Por otro lado, un debate más cercano a la práctica española: al tratarse de una deducción en cuota con base máxima elevada y dirigida por definición a contribuyentes con capacidad de inmovilizar capital en inversiones de alto riesgo durante años, es un incentivo que solo puede aprovechar quien ya dispone de un patrimonio y una renta considerables, lo que sitúa a esta figura en el extremo opuesto del espectro distributivo que hemos ido trazando a lo largo del apartado, pese a operar formalmente en cuota y no en base.

\subsubsection{Mecenazgo y donativos}
Los gastos fiscales en virtud del mecenazgo y las donaciones altruistas buscan complementar la financiación pública de actividades de interés general con financiación privada canalizada a través de entidades sin ánimo de lucro.

En España, el art. 68.3 LIRPF articula la deducción por donativos con una estructura escalonada: $80\%$ sobre los primeros 250 euros donados, y $40\%$ sobre el exceso, porcentaje que se eleva al $45\%$ cuando el contribuyente ha donado a la misma entidad, por un importe igual o superior, en cada uno de los dos periodos impositivos inmediatamente anteriores, sirviendo como una prima a la fidelidad del donante que busca fomentar no solo la donación puntual sino el compromiso sostenido con una causa.

El debate académico sobre los incentivos fiscales al mecenazgo tiene su origen en un trabajo pionero de FELDSTEIN (1975), que introdujo el concepto de elasticidad-precio de la donación: si un euro donado cuesta al contribuyente menos de un euro real por efecto de la deducción fiscal (en el ejemplo español, entre 0,20 y 0,60 euros, según el tramo), cabe preguntarse cuánto varía el volumen de donaciones ante cambios en ese precio fiscal de donar. La literatura posterior, sistematizada por autores como ANDREONI, distingue entre donantes cuya generosidad es sensible al incentivo fiscal y donantes que habrían donado igualmente sin el incentivo. Es decir, el mismo problema conceptual de peso muerto que vimos en el bloque de previsión social, aplicado ahora a la filantropía.

Este debate no está cerrado ni en España ni en la literatura internacional, y explica por qué las reformas del régimen de mecenazgo suelen justificarse invocando el argumento de abaratar la donación para maximizar su volumen, sin que exista, hasta la fecha, una evaluación de impacto pública y sistemática para el caso español equivalente a la disponible en otros países.

\subsubsection{Movilidad internacional del trabajo y empleo}
Las exenciones previstas en virtud de la movilidad internacional del trabajo buscan facilitar la proyección internacional de trabajadores y empresas españolas, en un contexto de internacionalización creciente de la economía desde finales de los noventa.

En España, el art. 7.p) LIRPF exime los rendimientos del trabajo obtenidos por trabajos efectivamente realizados en el extranjero para una empresa o entidad no residente, con el límite de $60.100$ euros anuales, siempre que en el país de destino exista un impuesto de naturaleza análoga al IRPF y no se trate de un paraíso fiscal. 

Es una de las exenciones más antiguas del catálogo introducida ya por la Ley 40/1998 con un límite inicial de 21.000 euros, elevado en el 2000 a 60.010 euros y prácticamente inalterado desde entonces: un cuarto de siglo de vigencia casi sin retoques cuantitativos, en fuerte contraste con la volatilidad normativa que hemos visto en otros bloques de este catálogo. Junto a esta exención, procede recordar la exención de indemnizaciones por despido (límite de 180.000 euros) como beneficio fiscal vinculado a la protección del empleo, y las rentas en especie exentas del art. 42.3 LIRPF (seguro de enfermedad, cheques-restaurante, transporte) como incentivos indirectos a la retribución flexible en el seno de la relación laboral.

El contraste entre este régimen de expatriados (art. 7.p) y el régimen de impatriados que vimos antes (Ley Beckham) es un excelente cierre: España ha construido, a lo largo de un cuarto de siglo, dos incentivos fiscales especulares que comparten la misma lógica de fondo (movilidad internacional del capital humano como bien a incentivar) pero que revelan objetivos de política económica distintos según el sentido del flujo, y que rara vez se estudian juntos pese a iluminarse mutuamente.

\subsection{Dispersión territorial: gastos fiscales autonómicos y forales}
Ya establecimos, al fijar el marco competencial en el apartado 1, que la Ley 22/2009 cede a las Comunidades Autónomas de régimen común capacidad normativa parcial sobre el IRPF, incluida —conforme a su art. 46— la aprobación de deducciones propias sobre la cuota íntegra autonómica en tres ámbitos: circunstancias personales y familiares, inversiones no empresariales, y aplicación de renta. 

Lo que aquí toca no es repetir esa competencia, ya fijada, sino mostrar su materialización concreta: qué ha hecho cada Comunidad con ese margen normativo, y qué problemas genera esa materialización.

\subsubsection{Magnitud del fenómeno: un catálogo que se ha vuelto inabarcable}
Para la campaña de la Renta de 2025, el conjunto de las Comunidades Autónomas de régimen común tenía vigentes aproximadamente $336$ deducciones autonómicas distintas, una cifra que por sí sola ilustra el problema mejor que cualquier explicación abstracta. 

Y no se trata solo de un problema de volumen, sino de heterogeneidad de diseño dentro de cada categoría temática: dentro del bloque familia que vimos, cada Comunidad ha desarrollado variantes propias con requisitos, umbrales de renta y cuantías completamente distintos entre sí, de modo que dos familias con idéntica composición y renta, residentes en Comunidades distintas, pueden terminar con resultados fiscales sustancialmente diferentes sin que exista ninguna diferencia relevante en su capacidad económica real.

Para dar una idea concreta de hasta qué punto de detalle ha llegado esta dispersión, algunos ejemplos recientes son ilustrativos por su singularidad: La Rioja ha introducido una deducción específica de $250$ euros para contribuyentes con enfermedad celíaca diagnosticada (por el sobrecoste de los productos sin gluten); Galicia ha aprobado deducciones para neutralizar la tributación de indemnizaciones a afectados por ELA y por el síndrome de la talidomida, o una nueva deducción del $15\%$ por libros de texto y material escolar con un límite de 105 euros; y varias Comunidades han desarrollado deducciones propias por instalación de vehículos eléctricos o puntos de recarga, o por asentamiento en municipios en riesgo de despoblación. Esta última categoría, de hecho, revela una finalidad de política pública que no aparece en absoluto en el catálogo estatal: precisamente porque el problema de la despoblación rural afecta de forma muy desigual al territorio español y solo tiene sentido como política diferenciada por Comunidad.

\subsubsection{Implicaciones: ¿equidad?}
El primer problema que plantea esta dispersión es de \textbf{equidad horizontal territorial}: el art. 139.1 CE proclama que todos los españoles tienen los mismos derechos y obligaciones en cualquier parte del territorio del Estado, principio que convive en tensión permanente con el art. 156 CE, que reconoce a las Comunidades Autónomas autonomía financiera para el desarrollo y ejecución de sus competencias. El Tribunal Constitucional ha resuelto reiteradamente esta tensión a favor de la autonomía financiera, admitiendo diferencias de trato tributario entre Comunidades como consecuencia lógica y querida del propio modelo de Estado autonómico, siempre que no se trate de discriminaciones arbitrarias. Pero la existencia de una doctrina constitucional que la avala no elimina el problema de fondo que late en cualquier análisis riguroso de gasto fiscal: dos contribuyentes con idéntica capacidad económica, separados solo por una frontera autonómica, pueden terminar soportando cargas tributarias sustancialmente distintas, sin que exista diferencia alguna en el hecho imponible ni en la base que las justifique.

El segundo problema es el de la \textbf{competencia fiscal} entre Comunidades, un debate mucho más desarrollado en la literatura y en el debate político español a propósito del Impuesto sobre Sucesiones y Donaciones, pero que empieza a reproducirse también en el terreno de las deducciones del IRPF, especialmente en las dirigidas a atraer residentes de renta alta o inversión empresarial. La lógica de fondo es la misma que analizamos con SHS y la movilidad internacional del capital al hablar del modelo dual: si los contribuyentes de renta alta pueden elegir su Comunidad de residencia atendiendo, entre otros factores, a la fiscalidad, las Comunidades Autónomas compiten entre sí por atraerlos, con el riesgo de una carrera a la baja que erosiona la capacidad recaudatoria conjunta del sistema, sin que exista, a diferencia del ámbito internacional, un marco de coordinación fiscal vinculante equivalente al que intentan construir organismos como la OCDE para la competencia entre Estados.

El tercer problema, más práctico pero no menos relevante, es el de la \textbf{infrautilización efectiva} de este catálogo por los propios contribuyentes a quienes se dirige. Al no figurar de forma automática en el borrador que confecciona la AEAT, que solo integra con seguridad los datos de los que dispone por suministro de terceros, buena parte de las deducciones autonómicas exigen que el propio contribuyente las identifique y las marque activamente, con la consecuencia documentada de que un porcentaje significativo de los beneficiarios potenciales simplemente no las aplica por desconocimiento, en una estimación agregada reciente que sitúa la pérdida colectiva de estos derechos no ejercidos en \textbf{varios miles de millones de euros anuales} en el conjunto del territorio nacional. Este dato conecta directamente con el debate de invisibilidad presupuestaria que cerramos en el subapartado 1: si allí el problema era que el Parlamento no escrutina el gasto fiscal con la misma intensidad que el gasto directo, aquí el problema es su reverso desde la perspectiva del contribuyente: un beneficio fiscal que exige activación consciente del beneficiario, tiende sistemáticamente a infrautilizarse entre quienes menos familiarizados están con el sistema tributario, que no son necesariamente quienes menos lo necesitan.

\subsection{Balance: magnitud, eficacia y debate de reforma}
Cerramos el apartado con la pregunta que ha quedado pendiente: ¿funcionan realmente los gastos fiscales que hemos catalogado? No es una pregunta retórica. Existe, en España, un cuerpo de evaluación empírica reciente y accesible que permite responderla con datos, y es el que desarrollo aquí.

\subsubsection{Magnitud cuantitativa}
Según los datos del Instituto de Estudios Fiscales, el conjunto de beneficios fiscales de todo el sistema tributario español se ha situado, en la serie 2010-2022, en una horquilla de entre el $5\%$ y el $8,3\%$ del PIB, con una tendencia descendente desde el máximo del periodo ($8,3\%$ en 2010, en plena crisis, cuando la recaudación caía más deprisa que el coste nominal de los beneficios fiscales) hasta estabilizarse en torno al $6\%$ en los últimos ejercicios de la serie.

Centrado ya en el Presupuesto de Beneficios Fiscales del Estado, la serie reciente permite una lectura de conjunto con cifras concretas: $41.939$ millones de euros en el proyecto de Presupuestos para 2022 (un incremento del $9,7\%$ sobre el ejercicio anterior), de los que el IRPF representó $11.221$ millones (el $27\%$ del total); y $45.269$ millones de euros en el proyecto para 2023, de los que el IRPF representó $11.179$ millones (en torno al $25\%$ del total, con un descenso interanual del $1,6\%$.º\footnote{Explicado, entre otros factores, por la caída de $190$ millones en la reducción por aportaciones a sistemas de previsión social (resultado directo del recorte del límite de aportación deducible) y de $98$ millones en la deducción transitoria por vivienda habitual, que sigue reduciéndose año a año a medida que se amortizan las hipotecas cubiertas por el régimen transitorio de 2013
}

Puesto en perspectiva, un coste anual del orden de $11.000$ millones de euros en el IRPF representa una cifra comparable a una parte sustancial de la recaudación líquida anual del propio impuesto, lo que da la medida real de por qué no es una cuestión menor del tema, sino una pieza de magnitud presupuestaria equivalente a la del propio diseño estructural que analizamos.

\subsubsection{Evaluación de eficacia: Informe de Gasto de la AIReF}
El origen de la evaluación sistemática de beneficios fiscales en España es reciente y tiene un hito preciso: por encargo del Consejo de Ministros, la Autoridad Independiente de Responsabilidad Fiscal (AIReF) incorporó los beneficios fiscales como uno de los cuatro bloques de su Informe de Gasto, junto con el gasto hospitalario del SNS, los incentivos a la contratación y las infraestructuras de transporte. 

El primer estudio de esta fase, presentado por la AIReF en 2020, analizó $13$ beneficios fiscales con un coste conjunto de $35.000$ millones de euros, equivalente al $60\%$ del total de beneficios fiscales existentes en el sistema tributario español. Para cada uno de los trece, la AIReF aplicó una metodología homogénea que evalúa tres preguntas: \textbf{si el beneficio cumple el objetivo para el que fue creado, si genera alguna distorsión o ineficiencia adicional, y si constituye la mejor vía posible para alcanzar ese objetivo frente a alternativas de gasto directo}.\footnote{Este primer estudio de la AIReF tuvo continuidad institucional: el Ministerio de Hacienda y Función Pública, encargó a la unidad de evaluación del propio Instituto de Estudios Fiscales la evaluación de otros $15$ beneficios fiscales adicionales durante 2021-2022, aplicando la misma metodología que la AIReF, con resultados publicados en el \textit{Informe de Revisión de Beneficios Fiscales de 2021} (2022).
}

Este seguimiento institucionalizado es, en sustancia, la aplicación práctica española de la misma filosofía de rendición de cuentas que motivó la reforma presupuestaria estadounidense de 1974: convertir el beneficio fiscal, hasta ahora invisible al escrutinio ordinario, en objeto de evaluación sistemática y periódica.

Dos resultados concretos de esa evaluación conectan directamente con contenido que ya hemos desarrollado en este apartado, y conviene traerlos aquí como prueba empírica, no como intuición doctrinal.

\paragraph{Resultado 1: Reducción de aportaciones a Sistemas de Previsión Social}
La Autoridad concluye expresamente que el beneficio fiscal no consigue el objetivo de incentivar el ahorro a largo plazo, y plantea su reformulación completa, en línea con las recomendaciones del Pacto de Toledo sobre ahorro complementario. Es la confirmación empírica, con metodología institucional española y no solo con el precedente estadounidense de DUFLO et al. (2006) que ya citamos, de que el diseño como reducción en base es un problema de eficacia demostrado y no solo teórico (regresivo y de dudoso efecto sobre el comportamiento de ahorro real, frente a la simple sustitución de ahorro que ya se iba a producir).

\paragraph{Resultado 2: Tributación conjunto}
La tributación conjunta, por el contrario, sí supera la evaluación: la AIReF concluye que el instrumento cumple razonablemente su objetivo declarado de adecuar el impuesto a la composición de rentas del hogar. Este resultado es importante para matizar que no todo gasto fiscal evaluado resulta ineficaz, y conviene evitar la simplificación de que toda desviación del benchmark es, por definición, una mala política.
 

\subsubsection{Debate de reforma: ensanchar la base, evaluar antes de crear}
La doctrina de referencia en el ámbito comparado ya ha sido analizada: ensanchar la base impositiva eliminando beneficios fiscales de eficacia dudosa, para poder sostener tipos marginales más bajos con la misma recaudación.

[Probable] La corriente de reforma que se deriva de esta evidencia, y que ya empieza a materializarse institucionalmente en España con el ciclo del Informe de Gasto 2022-2026 anunciado por la AIReF, apunta en una dirección clara: someter cualquier beneficio fiscal nuevo a evaluación ex ante antes de su aprobación, y revisar periódicamente los ya existentes (con cláusulas de caducidad o, como mínimo, con obligación de reevaluación), en lugar de dejar que se acumulen indefinidamente por el problema de invisibilidad presupuestaria. 

Es, en última instancia, la traducción institucional española de la intuición original de SURREY (1973): si un beneficio fiscal es, en sustancia, gasto público, debe estar sometido al mismo estándar de evaluación y rendición de cuentas que cualquier programa de gasto directo.






\clearpage
\section{Conclusión} 


\subsection{Recapitulación}


\subsection{Extensiones y relación con otras partes del Temario}



\subsection{Opinión}

\subsubsection{Idea final (Salida o cierra)}

\clearpage
\pagenumbering{gobble}
\MyBibliography


Ley 44/1978, de 8 de septiembre, del Impuesto sobre la Renta de las Personas Físicas.
STC 45/1989, de 20 de febrero (BOE de 21 de febrero de 1989).
Ley 20/1989, de 28 de julio, de adaptación del IRPF y del Impuesto Extraordinario sobre el Patrimonio.
Ley 18/1991, de 6 de junio, del Impuesto sobre la Renta de las Personas Físicas.
Ley 40/1998, de 9 de diciembre, del Impuesto sobre la Renta de las Personas Físicas y otras Normas Tributarias.
Real Decreto Legislativo 3/2004, de 5 de marzo, Texto Refundido de la Ley del IRPF.
Ley 35/2006, de 28 de noviembre, del Impuesto sobre la Renta de las Personas Físicas.
Ley 22/2009, de 18 de diciembre, por la que se regula el sistema de financiación de las Comunidades Autónomas y se modifican determinadas normas tributarias.
Ley 26/2014, de 27 de noviembre, por la que se modifican la Ley 35/2006 del IRPF, el TRLIRNR y otras normas tributarias.
Ley 11/2020, de 30 de diciembre, de Presupuestos Generales del Estado para 2021.
Ley 7/2024, de 20 de diciembre, por la que se establecen un Impuesto Complementario para garantizar un nivel mínimo global de imposición y se modifican otras normas tributarias.
Comisión de Expertos para la Reforma del Sistema Tributario Español (2014): Informe (Informe Lagares II), Ministerio de Hacienda.
Ministerio de Hacienda (1998): Informe de la Comisión para el Estudio y Propuesta de Medidas para la Reforma del IRPF (Informe Lagares I, 1997-1998).
ZORNOZA PÉREZ, J. (1989): "Aspectos constitucionales del régimen de tributación conjunta en el IRPF (Comentario a la STC 45/1989)", Revista Española de Derecho Constitucional.
PALAO TABOADA, C. (1981): "El tratamiento de la familia en la imposición sobre la renta", Civitas, Revista Española de Derecho Financiero.
Real Decreto 439/2007, de 30 de marzo, por el que se aprueba el Reglamento del IRPF.
Ley 58/2003, de 17 de diciembre, General Tributaria.
Ley Orgánica 8/1980, de 22 de septiembre, de Financiación de las Comunidades Autónomas (LOFCA).
Ley 12/2002, de 23 de mayo, por la que se aprueba el Concierto Económico con el País Vasco.
Ley 28/1990, de 26 de diciembre, por la que se aprueba el Convenio Económico entre el Estado y la Comunidad Foral de Navarra.
Ley 28/2022, de 21 de diciembre, de fomento del ecosistema de las empresas emergentes ("Ley de Startups").
ALMUDÍ CID, J. y SERRANO ANTÓN, F.: "La residencia fiscal de las personas físicas en los convenios de doble imposición internacional y en la normativa interna española", Estudios Financieros, CEF.
Constitución Española de 1978, art. 134.2.
Ley 47/2003, de 26 de noviembre, General Presupuestaria (art. 33.2 y 37.2).
Ley 41/1994, de 30 de diciembre, de Presupuestos Generales del Estado para 1995 (disposición adicional 24ª).
SURREY, S. S. (1973): Pathways to Tax Reform: The Concept of Tax Expenditures*, Harvard University Press.
SURREY, S. S. y McDANIEL, P. R. (1985): *Tax Expenditures*, Harvard University Press.
Ministerio de Hacienda: *Memoria de Beneficios Fiscales*, Presupuestos Generales del Estado (edición anual).
DOMÍNGUEZ MARTÍNEZ, J. M. (2018): "El concepto de renta fiscal según Haig-Simons: los vasos comunicantes", *eXtoikos*, nº 21.
SCHANZ, G. von (1896): "Der Einkommensbegriff und die Einkommensteuergesetze", *FinanzArchiv*, 13.
Ley 11/2021, de 9 de julio, de medidas de prevención y lucha contra el fraude fiscal (modificación del art. 36 LIRPF sobre pactos sucesorios).
STS 407/2016, de 9 de febrero (Sala de lo Contencioso-Administrativo), sobre la calificación mortis causa de los pactos sucesorios a efectos del art. 33.3.b) LIRPF.
RINCÓN DE PABLO, G. (2014): "La lenta agonía de los coeficientes de abatimiento y la muerte de los coeficientes de corrección monetaria", Observatorio Inmobiliario, Garrigues.
SLEMROD, J. y CHEN, X. (2023): "Are capital gains the Achilles' heel of taxing the rich?", Oxford Review of Economic Policy, vol. 39(3).
Ley 35/2006, arts. 42, 43, 85 (ya citada; ahora en su contenido específico).
Ley 26/2014 (ya citada; introducción del art. 27.1, tercer párrafo LIRPF).
STS 1443/2025, de 12 de noviembre (Sala de lo Contencioso-Administrativo), sobre calificación de rendimientos de socios profesionales.
SØRENSEN, P. B. (2005): "Neutral Taxation of Shareholder Income", International Tax and Public Finance.
CNOSSEN, S. (2000): "Taxing Capital Income in the Nordic Countries: A Model for the European Union?", FinanzArchiv.
MIRRLEES, J. et al. (2011): Tax by Design: The Mirrlees Review, Institute for Fiscal Studies / Oxford University Press.
Directiva 88/361/CEE del Consejo, de 24 de junio de 1988, para la aplicación del artículo 67 del Tratado (liberalización de movimientos de capitales).
BRYS, B. (2014): "Trends and Reform Options for the Taxation of Capital Income in OECD Countries", OECD Taxation Working Papers.
DURÁN CABRÉ, J. M. (2001): "El Impuesto Dual. Estudio teórico y análisis empírico para el caso español", VIII Encuentros en Economía Pública, Cáceres.
DÍAZ CARO, C., ONRUBIA FERNÁNDEZ, J. y PÉREZ MAYO, J. (2013): "Progresividad y redistribución por fuentes de renta en el IRPF dual", Hacienda Pública Española / Review of Public Economics, nº 206.
SØRENSEN, P. B. (1994): "From the Global Income Tax to the Dual Income Tax: Recent Tax Reforms in the Nordic Countries", International Tax and Public Finance.
Ley 35/2006, arts. 81 y 81 bis (ya citada; ahora en su contenido específico).
Real Decreto-ley 1/2015, de 27 de febrero, de mecanismo de segunda oportunidad, reducción de carga financiera y otras medidas de orden social (introducción del art. 81 bis LIRPF).
FRIEDMAN, M. (1962): Capitalism and Freedom, University of Chicago Press (formulación del impuesto negativo sobre la renta).
Ley 58/2003, General Tributaria, art. 8.d) (reserva de ley para beneficios fiscales).
Constitución Española de 1978, art. 31.1 (ya citado; ahora en su función de fundamento de la progresividad frente a la erosión por beneficios fiscales).
Congressional Budget and Impoundment Control Act of 1974 (Pub. L. 93-344), Estados Unidos.
HOWARD, C. (1997): The Hidden Welfare State: Tax Expenditures and Social Policy in the United States, Princeton University Press.
Orden HFP/105/2017, de 6 de febrero (regularización de cobros indebidos del abono anticipado del art. 81 bis LIRPF mediante Modelo 122).
Ley 16/2012, de 27 de diciembre (ya citada; ahora en su contenido sobre derogación de la deducción por vivienda).
Ley 12/2023, de 24 de mayo, por el derecho a la vivienda.
Real Decreto-ley 19/2021, de 5 de octubre, de medidas urgentes para impulsar la actividad de rehabilitación edificatoria.
Real Decreto-ley 20/2020, de 29 de mayo, por el que se establece el Ingreso Mínimo Vital.
Ley 49/2002, de 23 de diciembre, de régimen fiscal de las entidades sin fines lucrativos y de los incentivos fiscales al mecenazgo.
Real Decreto-ley 3/2000, de 23 de junio (elevación del límite del art. 7.p LIRPF).
POTERBA, J. M. (1984): "Tax Subsidies to Owner-Occupied Housing: An Asset-Market Approach", Quarterly Journal of Economics, vol. 99.
FELDSTEIN, M. (1975): "The Income Tax and Charitable Contributions: Part I—Aggregate and Distributional Effects", National Tax Journal, vol. 28.
DUFLO, E., GALE, W., LIEBMAN, J., ORSZAG, P. y SAEZ, E. (2006): "Saving Incentives for Low- and Middle-Income Families: Evidence from a Field Experiment with H\&R Block", Quarterly Journal of Economics, vol. 121.
KERR, W., LERNER, J. y SCHOAR, A. (2014): "The Consequences of Entrepreneurial Finance: Evidence from Angel Financings", Review of Financial Studies, vol. 27.
Ley 22/2009, art. 46 (ya citada; ahora en su contenido específico sobre competencias normativas en deducciones autonómicas).
Concierto Económico con el País Vasco (Ley 12/2002, ya citada), arts. 6 a 13 (regulación del IRPF concertado).
Ley 3/1989, de 30 de mayo, del Parlamento Vasco, de armonización, coordinación y colaboración fiscal entre las Instituciones de los Territorios Históricos.
Constitución Española de 1978, arts. 139.1 y 156 (ya citado este último de forma indirecta; ahora en su función de fundamento de la autonomía financiera autonómica).
AIReF (2020): *Spending Review, Fase II: Beneficios Fiscales*, Autoridad Independiente de Responsabilidad Fiscal.
Ministerio de Hacienda y Función Pública (2022): *Informe de Revisión de Beneficios Fiscales 2021*, marzo de 2022.
GARCÍA-HERRERA, C. (2023): "Evaluación de beneficios fiscales: la experiencia española", Seminario de Política Fiscal, CEPAL, Santiago de Chile.
OCDE: *Tax Expenditures in OECD Countries*, informes periódicos de la Organización para la Cooperación y el Desarrollo Económicos.
Proyectos de Ley de Presupuestos Generales del Estado 2022 y 2023, Memoria de Beneficios Fiscales (ya citada de forma genérica en el apartado 2; ahora con cifras concretas de ambos ejercicios).

\MyBibItem{DAWID y GATTI}{Chapter 2 - Agent-Based Macroeconomics}{2018}{https://doi.org/10.1016/bs.hescom.2018.02.006}


% ANEXOS
\clearpage










\end{document}